% ==============================================================================
% CH13_AI_SYSTEM_DESIGN.TEX — CHƯƠNG 13: THIẾT KẾ KIẾN TRÚC HỆ THỐNG AI QUY MÔ LỚN
% ==============================================================================
% Định dạng: Sách chuyên khảo kinh điển (Classic Academic & Engineering Monograph)
% Chuẩn mực: Tuyệt đối không dùng khung hộp 4 phía, toán học 5 bước, bài tập đa tầng
% ==============================================================================

\chapter{Thiết Kế Kiến Trúc Hệ Thống AI Quy Mô Lớn}
\label{ch:ai-system-design}

Trong quá trình phỏng vấn kỹ thuật cấp cao (Staff / Principal AI Engineer) cũng như khi thiết kế hạ tầng thực tế cho doanh nghiệp, năng lực phân biệt một lập trình viên thông thường với một kiến trúc sư hệ thống AI (\textit{AI System Architect}) nằm ở khả năng: \textbf{Thiết kế một hệ thống AI phục vụ hàng chục triệu người dùng đồng thời, đảm bảo độ trễ P99 dưới 1 giây, tính sẵn sàng 99.99\%, và tối ưu hóa hàng trăm ngàn USD chi phí phần cứng GPU mỗi tháng.}

Thiết kế hệ thống AI (\textit{AI System Design}) không đơn thuần là thiết kế web backend truyền thống (CRUD, SQL, Load Balancer). Các hệ thống AI sở hữu những đặc thù phần cứng và tính toán hoàn toàn khác biệt:
\begin{itemize}
  \item \textbf{Cực kỳ ngốn tài nguyên tính toán (Compute-Bound) \& Bộ nhớ (Memory-Bound)}: Mỗi lượt sinh token đòi hỏi nạp hàng chục Gigabyte trọng số qua băng thông VRAM của GPU.
  \item \textbf{Độ trễ không tất định (Non-Deterministic Latency)}: Một truy vấn trả về 10 tokens chỉ mất $200\text{ms}$, nhưng một truy vấn sinh 2000 tokens có thể kéo dài tới $30\text{ giây}$.
  \item \textbf{Phân mảnh bộ nhớ động (KV Cache Allocation)}: Cần quản lý hàng triệu block bộ nhớ trên GPU tương tự cơ chế bộ nhớ ảo của hệ điều hành.
\end{itemize}

Chương này trang bị cho bạn khung phương pháp chuẩn mực (\textit{System Design Framework}), kỹ năng ước tính dung lượng thô (\textit{Back-of-the-Envelope Estimation}), kiến trúc bộ nhớ đệm đa tầng (\textit{Multi-Tier Semantic Caching}), cơ chế định tuyến phân cấp (\textit{Hierarchical Cascading Routing}), và giải thuật cân bằng tải tối ưu cho cụm máy chủ vLLM.

---

\section{Khung Phương Pháp Thiết Kế Hệ Thống AI Chuẩn Mực}
\label{sec:ch13-framework}

Để giải quyết một bài toán thiết kế hệ thống mở mà không bị sa đà vào các chi tiết vụn vặt, kỹ sư AI luôn áp dụng quy trình 4 bước có cấu trúc:

\begin{figure}[htbp]
\centering
\small
\begin{verbatim}
+------------------------------------------------------------------------+
| BƯỚC 1: LÀM RÕ YÊU CẦU (Requirements Clarification)                   |
| * Chức năng: Hỏi đáp RAG, Code generation, Agent tự hành              |
| * Phi chức năng: QPS đỉnh, TTFT, TPOT, P99 Latency, SLA 99.99%, Budget  |
+------------------------------------------------------------------------+
                                   |
                                   v
+------------------------------------------------------------------------+
| BƯỚC 2: ƯỚC TÍNH DUNG LƯỢNG THÔ (Back-of-the-Envelope Estimation)      |
| * QPS trung bình / đỉnh, Tokens/s, Băng thông mạng, VRAM KV Cache      |
| * Số lượng máy chủ GPU cần thiết, Phân tích chi phí FinOps             |
+------------------------------------------------------------------------+
                                   |
                                   v
+------------------------------------------------------------------------+
| BƯỚC 3: THIẾT KẾ KIẾN TRÚC TỔNG THỂ (High-Level Architecture)          |
| * API Gateway, Multi-Tier Caching (Exact + Semantic), Model Router     |
| * Cụm Inference vLLM / Triton, Vector DB Sharding, RAG Pipeline        |
+------------------------------------------------------------------------+
                                   |
                                   v
+------------------------------------------------------------------------+
| BƯỚC 4: ĐI SÂU VÀO ĐIỂM NGHẼN & TỐI ƯU HÓA (Deep-Dive & Bottlenecks)   |
| * Prefix Caching Affinity, Least-Loaded KV Balancing, Cascading Storm  |
| * Rate Limiting phân tán, Circuit Breaker ngân sách, Graceful Fallback |
+------------------------------------------------------------------------+
\end{verbatim}
\caption{Quy trình 4 bước chuẩn mực trong Thiết kế Hệ thống AI Doanh nghiệp}
\label{fig:ai-sys-design-framework}
\end{figure}

\subsection{Các Chỉ Số Hiệu Năng Cốt Lõi (SLA / SLI Metrics)}

Khi bàn về hiệu năng hệ thống AI, tuyệt đối không dùng từ ``nhanh'' hay ``chậm'' chung chung. Bắt buộc phải sử dụng các đại lượng kỹ thuật chuẩn xác:
\begin{itemize}
  \item \textbf{TTFT (Time To First Token)}: Thời gian từ khi Gateway nhận được request của người dùng cho đến khi token đầu tiên được sinh ra và gửi về client. Phản ánh thời gian thực thi giai đoạn \textit{Prefill} của mô hình cộng với độ trễ hàng đợi mạng. Tiêu chuẩn vàng: $\text{TTFT} \le 500\text{ms}$.
  \item \textbf{TPOT (Time Per Output Token)}: Thời gian để mô hình sinh ra một token tiếp theo trong giai đoạn \textit{Decode}. Quyết định tốc độ đọc của người dùng. Tốc độ đọc tự nhiên của mắt người là khoảng $15 - 20\text{ tokens/giây}$ (tương ứng $\text{TPOT} \approx 50 - 65\text{ms/token}$).
  \item \textbf{P99 Latency (Percentile 99)}: Thời gian phản hồi mà $99\%$ người dùng trải nghiệm thấp hơn con số này. Đo lường mức độ ổn định của hệ thống trước hiện tượng nghẽn hàng đợi (\textit{Head-of-Line Blocking}).
  \item \textbf{Throughput (Tokens per Second)}: Tổng số lượng token mà toàn bộ cụm máy chủ GPU có thể xử lý và sinh ra trong một giây.
\end{itemize}

---

\section{Ước Tính Dung Lượng Thô: Bài Toán 10 Triệu Người Dùng (10M DAU)}
\label{sec:ch13-estimation}

Hãy cùng thực hiện một bài toán ước tính dung lượng kinh điển trong phỏng vấn: \textbf{Thiết kế hệ thống Trợ lý Hỏi Đáp Doanh nghiệp phục vụ 10 triệu người dùng hoạt động hàng ngày (10M DAU).}

\subsection{1. Ước Tính Số Lượng Truy Vấn \& QPS Đỉnh}
\begin{itemize}
  \item Số người dùng hoạt động hàng ngày: $\text{DAU} = 10.000.000$ users.
  \item Mỗi người dùng thực hiện trung bình $10$ truy vấn mỗi ngày.
  \item Tổng số truy vấn mỗi ngày:
  \begin{equation}
  N_{\text{queries}} = 10.000.000 \times 10 = 100.000.000 \text{ requests/ngày}
  \end{equation}
  \item QPS trung bình trong một ngày ($86.400$ giây):
  \begin{equation}
  \text{QPS}_{\text{avg}} = \frac{100.000.000}{86.400} \approx 1.157 \text{ queries/giây}
  \end{equation}
  \item Lưu lượng thực tế không phân bổ đều mà có các khung giờ cao điểm (Peak Hours). Áp dụng hệ số tải đỉnh (\textit{Peak Factor}) $k_{\text{peak}} = 3$:
  \begin{equation}
  \text{QPS}_{\text{peak}} = 1.157 \times 3 \approx \mathbf{3.500 \text{ queries/giây}}
  \end{equation}
\end{itemize}

\subsection{2. Ước Tính Băng Thông Mạng (Network Bandwidth)}
Giả định thông số độ dài văn bản:
\begin{itemize}
  \item Prompt đầu vào (bao gồm câu hỏi người dùng + tài liệu ngữ cảnh RAG): Trung bình $L_{\text{in}} = 1.000$ tokens.
  \item Phản hồi đầu ra do mô hình sinh: Trung bình $L_{\text{out}} = 200$ tokens.
  \item Mỗi token tương đương trung bình $4$ bytes văn bản UTF-8.
\end{itemize}
Tổng lượng token mà cụm hệ thống cần luân chuyển mỗi giây vào giờ cao điểm:
\begin{equation}
\text{Tokens/s}_{\text{total}} = 3.500 \times (1.000 + 200) = 4.200.000 \text{ tokens/giây}
\end{equation}
Băng thông mạng tại tầng Gateway:
\begin{equation}
\text{Bandwidth} = 4.200.000 \times 4 \text{ bytes} \approx 16.8 \text{ MB/s} = \mathbf{134.4 \text{ Mbps}}
\end{equation}
Kết luận: Băng thông mạng ở tầng Gateway hoàn toàn nằm trong khả năng xử lý của các card mạng 1Gbps / 10Gbps thông thường. Điểm nghẽn không nằm ở mạng ngoài!

\subsection{3. Ước Tính Nhu Cầu Phần Cứng GPU \& Khủng Hoảng Chi Phí}
Điểm nghẽn thực sự nằm ở năng lực sinh token của cụm máy chủ GPU.
Tổng số lượng token cần \textbf{decode} mỗi giây vào giờ cao điểm:
\begin{equation}
\text{Decode Tokens/s} = 3.500 \text{ QPS} \times 200 \text{ tokens} = \mathbf{700.000 \text{ tokens/giây}}
\end{equation}

Giả sử doanh nghiệp triển khai mô hình mã nguồn mở Llama-3-70B ở định dạng FP8 trên máy chủ trang bị GPU NVIDIA H100 SXM5 (80GB HBM3).
Một nút máy chủ 4xH100 chạy vLLM với Tensor Parallelism ($TP=4$) và Continuous Batching có khả năng sinh thông lượng trung bình khoảng $2.500$ output tokens/giây.
Số lượng GPU H100 tối thiểu cần để gánh tải đỉnh mà không bị sập:
\begin{equation}
N_{\text{GPU}} = \frac{700.000 \text{ tokens/s}}{2.500 \text{ tokens/s/node}} \times 4 \text{ GPUs/node} = \mathbf{1.120 \text{ GPUs H100}}!
\end{equation}

\paragraph{Khủng hoảng chi phí phần cứng nếu không có kiến trúc tối ưu}:
Chi phí thuê một GPU NVIDIA H100 trên hạ tầng đám mây (AWS, CoreWeave, Lambda Labs) hiện nay rơi vào khoảng $\$2.50\text{/GPU-giờ}$.
Chi phí vận hành hàng tháng cho cụm $1.120$ GPU:
\begin{equation*}
\text{Cost}_{\text{month}} = 1.120 \times \$2.50 \times 24 \text{ giờ} \times 30 \text{ ngày} = \mathbf{\$2.016.000\text{/tháng}} \quad (\text{Hơn 2 triệu USD!})
\end{equation*}

Một doanh nghiệp không thể tồn tại với chi phí này nếu không có chiến lược tối ưu hóa kiến trúc! Là kiến trúc sư hệ thống, nhiệm vụ của bạn là kéo giảm con số này xuống \textbf{dưới $\$200.000\text{/tháng}$ (tiết kiệm hơn $90\%$)} thông qua hai trụ cột:
\begin{enumerate}
  \item \textbf{Multi-Tier Semantic Caching}: Giảm $35 - 50\%$ lưu lượng truy vấn phải chạm tới GPU.
  \item \textbf{Hierarchical Model Routing}: Chuyển $70\%$ câu hỏi đơn giản sang các mô hình nhỏ (Llama-3-8B) chạy trên GPU L4/A10G rẻ hơn gấp 5 lần!
\end{enumerate}

---

\section{Kiến Trúc Bộ Nhớ Đệm Ngữ Nghĩa Đa Tầng (Multi-Tier Semantic Caching)}
\label{sec:ch13-semantic-caching}

Trong các hệ thống web truyền thống, caching dựa trên sự trùng khớp chính xác của khóa (\textit{Exact Key Matching}): Ví dụ mã băm $\text{MD5}(\text{URL})$ hoặc truy vấn SQL.
Tuy nhiên, trong ngôn ngữ tự nhiên:
\begin{center}
\textit{``Thủ đô của nước Pháp là gì?''} \quad và \quad \textit{``Cho tôi biết tên thủ đô nước Pháp?''}
\end{center}
về mặt ký tự hoàn toàn khác nhau ($\text{Hash}_1 \neq \text{Hash}_2$), nhưng về mặt ngữ nghĩa là đồng nhất $100\%$ và có cùng câu trả lời: ``Paris''.
Nếu dùng Exact Cache, cả hai câu hỏi đều phải gọi tới mô hình LLM. Bộ nhớ đệm ngữ nghĩa (\textit{Semantic Cache}) giải quyết bài toán này.

\begin{figure}[htbp]
\centering
\small
\begin{verbatim}
[User Query q]
       |
       v
+-----------------------------------------------------------+
| TẦNG 1: Local In-Memory LRU Cache (Gateway Memory, <1ms)  |
+-----------------------------------------------------------+
       | (Miss)
       v
+-----------------------------------------------------------+
| TẦNG 2: Distributed Exact Hash Cache (Redis Cluster, <3ms)|
|   Key = SHA256(tenant_id + normalize(q))                  |
+-----------------------------------------------------------+
       | (Miss)
       v
+-----------------------------------------------------------+
| TẦNG 3: Vector Semantic Cache (Qdrant / Milvus HNSW, <20ms|
|   Tìm kiếm Cosine k-NN có Partition Filter theo Tenant     |
|   Nếu max(Similarity) >= tau (0.88) --> HIT                |
+-----------------------------------------------------------+
       | (Miss)
       v
+-----------------------------------------------------------+
| TẦNG 4: vLLM Inference Serving với Prefix Caching         |
+-----------------------------------------------------------+
\end{verbatim}
\caption{Kiến trúc bộ nhớ đệm đa tầng kết hợp Exact Hash Cache và Vector Semantic Cache}
\label{fig:two-tier-cache}
\end{figure}

\subsection{Cơ Chế Hoạt Động Toán Học Của Semantic Cache}
\begin{enumerate}
  \item \textbf{Tầng 1 \& 2: Exact Hash Cache (Redis)}:
  Tính mã băm chuẩn hóa của câu truy vấn sau khi chuyển thành chữ thường và loại bỏ khoảng trắng thừa:
  \begin{equation}
  k = \text{SHA256}(\text{normalize}(\mathbf{q}))
  \end{equation}
  Truy vấn Redis bằng thao tác $\mathcal{O}(1)$. Nếu trúng cache (\textit{Hit}), trả lời ngay lập tức trong vòng $1 - 2\text{ms}$. Tầng này xử lý khoảng $15 - 20\%$ lưu lượng (các câu hỏi lặp lại hoàn toàn do cùng một nhóm người dùng hỏi).

  \item \textbf{Tầng 3: Vector Semantic Cache (Qdrant / Milvus)}:
  Nếu Tầng 2 trượt (\textit{Miss}), ta tính vector nhúng ngữ nghĩa của truy vấn bằng mô hình embedding nhỏ và nhanh (như \code{bge-small-en-v1.5}):
  \begin{equation}
  \mathbf{e}_q = f_{\text{embed}}(\mathbf{q}) \in \mathbb{R}^d
  \end{equation}
  Thực hiện tìm kiếm k-hàng xóm gần nhất (k-NN) trong kho lưu trữ các truy vấn đã được giải đáp trước đó $\mathcal{C} = \{(\mathbf{e}_i, \mathbf{a}_i)\}_{i=1}^M$:
  \begin{equation}
  \mathbf{e}^* = \arg\max_{\mathbf{e}_i \in \mathcal{C}} \cos(\mathbf{e}_q, \mathbf{e}_i) = \arg\max_{\mathbf{e}_i \in \mathcal{C}} \frac{\mathbf{e}_q \cdot \mathbf{e}_i}{\|\mathbf{e}_q\| \|\mathbf{e}_i\|}
  \end{equation}
  Gọi độ tương đồng lớn nhất tìm được là $S_{\max} = \cos(\mathbf{e}_q, \mathbf{e}^*)$.
  Ta đặt một ngưỡng tương đồng khắt khe $\tau_{\text{threshold}}$:
  \begin{equation}
  \text{Trạng thái Cache} = \begin{cases}
  \text{\textbf{HIT}} \implies \text{Trả về câu trả lời } \mathbf{a}^* & \text{nếu } S_{\max} \ge \tau_{\text{threshold}} \\
  \text{\textbf{MISS}} \implies \text{Chuyển tiếp tới LLM} & \text{nếu } S_{\max} < \tau_{\text{threshold}}
  \end{cases}
  \end{equation}
\end{enumerate}

\paragraph{Đánh đổi của ngưỡng tương đồng $\tau_{\text{threshold}}$}:
\begin{itemize}
  \item Nếu chọn $\tau = 0.80$ (quá thấp): Hệ thống đạt tỷ lệ Hit Rate rất cao ($60\%$), nhưng sinh ra hiện tượng \textbf{Cache Pollution}: Người dùng hỏi ``Làm thế nào để nạp tiền vào tài khoản?'' có thể nhận nhầm câu trả lời của ``Làm thế nào để hủy tài khoản?''. Đây là lỗi nghiêm trọng!
  \item Nếu chọn $\tau = 0.98$ (quá cao): Câu hỏi phải gần như giống hệt nhau mới kích hoạt được cache, tỷ lệ Hit Rate giảm về dưới $5\%$, mất tác dụng tiết kiệm tài nguyên.
  \item \textbf{Thực tế kỹ nghệ}: Ngưỡng tối ưu thường nằm trong khoảng $\mathbf{\tau \in [0.88, 0.92]}$. Đồng thời, ta áp dụng thêm một bộ kiểm định thực thể (Entity Extractor): Chỉ chấp nhận Hit nếu hai câu hỏi có cùng tập thực thể chính (cùng mã cổ phiếu, cùng tên sản phẩm).
\end{itemize}

---

\section{Định Tuyến Phân Cấp: Mô Hình Nhỏ (SLM) Kết Hợp Mô Hình Lớn (LLM)}
\label{sec:ch13-hierarchical-routing}

Không phải mọi truy vấn của người dùng đều cần đến một mô hình khổng lồ 70B hoặc 405B tham số.
Trong thực tế doanh nghiệp:
\begin{itemize}
  \item Khoảng $60 - 70\%$ truy vấn là các tác vụ đơn giản: Phân loại ý định, trích xuất thực thể, kiểm tra chính tả, hoặc tóm tắt một đoạn văn ngắn.
  \item Chỉ có $30 - 40\%$ truy vấn thực sự đòi hỏi khả năng lập luận phức tạp: Giải bài toán đại số, thiết kế kiến trúc phần mềm, phân tích hợp đồng kinh tế nhiều trang.
\end{itemize}

Chiến lược \textbf{Hierarchical Model Routing} sử dụng một bộ phân loại siêu nhẹ (Router Classifier) đặt tại Gateway:

\begin{figure}[htbp]
\centering
\small
\begin{verbatim}
[Yêu Cầu Người Dùng]
        |
        v
+-------------------------------+
| Fast Classifier Router (2ms)  |
+-------------------------------+
     /                     \
 (Độ phức tạp <= 3)    (Độ phức tạp > 3)
   /                         \
  v                           v
+---------------------+     +----------------------+
| Tầng 1: SLM (8B)    |     | Tầng 2: LLM (70B)    |
| (GPU L4/A10G, $0.5) |     | (GPU H100 SXM, $2.5) |
+---------------------+     +----------------------+
\end{verbatim}
\caption{Cơ chế định tuyến phân cấp theo độ phức tạp truy vấn}
\label{fig:model-routing}
\end{figure}

Bộ phân loại Router có thể được huấn luyện bằng mô hình phân loại văn bản nhỏ (như DeBERTa-v3 hoặc Logistic Regression trên vector embedding) với độ trễ suy luận chỉ $2 - 5\text{ms}$. Thang điểm độ phức tạp được chấm từ 1 đến 5 dựa trên các đặc trưng:
\begin{itemize}
  \item Độ dài câu hỏi và độ dài ngữ cảnh cần xử lý.
  \item Sự xuất hiện của các từ khóa lập luận (``chứng minh'', ``so sánh đa chiều'', ``tối ưu hóa mã nguồn'').
  \item Số lượng bước suy luận dự kiến cần thực hiện.
\end{itemize}

\paragraph{Cascading Fallback \& Speculative Execution}:
Khi một câu hỏi được chuyển cho mô hình nhỏ 8B, hệ thống đo lường độ tự tin dựa trên điểm số xác suất logarit (\textit{Average Logprob Confidence}):
\begin{equation}
\mathcal{C} = \exp\left( \frac{1}{K} \sum_{i=1}^K \log P(y_i \mid y_{<i}, x) \right)
\label{eq:logprob-confidence}
\end{equation}
Nếu $\mathcal{C} < \theta_{\text{confidence}}$ (ví dụ $\theta = 0.70$), câu trả lời của mô hình nhỏ bị loại bỏ và yêu cầu tự động leo thang (\textit{Cascade}) lên mô hình lớn 70B! Cơ chế này đảm bảo độ chính xác của hệ thống luôn tiệm cận mô hình lớn nhất mà vẫn tiết kiệm được $70\%$ chi phí.

---

\section{Cân Bằng Tải \& Khả Năng Chống Chịu Sự Cố Cho Cụm vLLM}
\label{sec:ch13-load-balancing}

\subsection{Tại Sao Cân Bằng Tải Round-Robin Truyền Thống Thất Bại?}
Trong kiến trúc Web truyền thống, thuật toán \textit{Round-Robin} phân phối đều tuần tự các request tới các máy chủ backend: Request 1 đến Server A, Request 2 đến Server B.
Tuy nhiên, khi áp dụng Round-Robin vào cụm máy chủ LLM (vLLM / Triton), hệ thống sẽ nhanh chóng rơi vào thảm họa:
\begin{itemize}
  \item Giả sử Server A nhận phải 5 request liên tiếp đều yêu cầu sinh văn bản dài 2048 tokens $\implies$ VRAM của Server A bị đầy, KV Cache bị phân mảnh, độ trễ tăng vọt lên 20 giây!
  \item Trong khi đó, Server B nhận được các request chỉ cần phân loại 5 tokens $\implies$ GPU của Server B rảnh rỗi $90\%$ thời gian!
\end{itemize}

\subsection{Thuật Toán Least-KV-Cache-Tokens Có Tính Ngữ Cảnh (Prefix Affinity)}
Để tối ưu hóa cụm vLLM, kỹ sư hệ thống sử dụng thuật toán \textbf{Context-Aware Least-Loaded Load Balancer}:
Bộ cân bằng tải duy trì liên tục thông tin trạng thái từ API giám sát của các nút vLLM:
\begin{equation}
\text{Tải của nút } i: \quad \mathcal{T}_i = \alpha \cdot \text{ActiveRequests}_i + \beta \cdot \text{KVCacheUsagePercentage}_i
\end{equation}
Đặc biệt, hệ thống áp dụng kỹ thuật \textbf{Prefix Caching Affinity}:
Nếu một nhóm request cùng sử dụng chung một System Prompt dài (ví dụ một tài liệu cẩm nang nội bộ công ty dài 8000 tokens), Gateway sẽ ưu tiên định tuyến tất cả các request này về cùng một máy chủ vLLM nơi ma trận KV Cache của System Prompt đó đã nằm sẵn trên VRAM! Máy chủ này không cần tính toán lại giai đoạn Prefill 8000 tokens, giúp giảm TTFT từ $1200\text{ms}$ xuống còn vỏn vẹn \textbf{40ms}!

---

\section{Cài Đặt Bộ Cổng Định Tuyến \& Semantic Cache Hoàn Chỉnh}
\label{sec:ch13-implementation}

Dưới đây là mã nguồn Python chuẩn mực cài đặt lớp \code{ProductionSemanticCacheGateway}, kết hợp Exact Hash Cache, Semantic Vector Cache có phân tách Tenant (Multi-Tenant Isolation), và bộ định tuyến phân cấp mô hình:

\begin{lstlisting}[language=Python]
import hashlib
import math
from typing import Dict, Any, List, Optional, Tuple

class ProductionSemanticCacheGateway:
    """Gateway dieu phoi cache hai tang va dinh tuyen mo hinh an toan da khach hang."""
    
    def __init__(self, similarity_threshold: float = 0.88):
        self.similarity_threshold = similarity_threshold
        # Tang 1: Exact Hash Cache (SHA256 -> Answer)
        self.exact_cache: Dict[str, str] = {}
        # Tang 2: Vector Semantic Cache [(Vector, Question, Answer, TenantID)]
        self.vector_cache: List[Tuple[List[float], str, str, str]] = []

    @staticmethod
    def _normalize_and_hash(text: str, tenant_id: str = "default") -> str:
        """Chuan hoa chuoi va tinh ma bam ket hop voi tenant_id."""
        cleaned = " ".join(text.strip().lower().split())
        compound = f"{tenant_id}::{cleaned}"
        return hashlib.sha256(compound.encode("utf-8")).hexdigest()

    @staticmethod
    def _cosine_similarity(v1: List[float], v2: List[float]) -> float:
        """Tinh do tuong dong Cosine giua 2 vector nhung."""
        dot = sum(a * b for a, b in zip(v1, v2))
        norm1 = math.sqrt(sum(a * a for a in v1))
        norm2 = math.sqrt(sum(b * b for b in v2))
        if norm1 == 0.0 or norm2 == 0.0:
            return 0.0
        return dot / (norm1 * norm2)

    def query(
        self, 
        prompt: str, 
        embedding: Optional[List[float]] = None,
        tenant_id: str = "default"
    ) -> Tuple[str, str, Dict[str, Any]]:
        """Xu ly truy van qua quy trinh 2 tang cache co phan tach tenant."""
        # 1. Kiem tra Tang 1: Exact Hash Cache
        hash_key = self._normalize_and_hash(prompt, tenant_id)
        if hash_key in self.exact_cache:
            return self.exact_cache[hash_key], "EXACT_CACHE_HIT", {"latency_ms": 1.5}

        # 2. Kiem tra Tang 2: Vector Semantic Cache
        if embedding is not None and len(self.vector_cache) > 0:
            best_sim = -1.0
            best_answer = ""
            matched_q = ""
            for cached_emb, cached_q, cached_a, cached_tenant in self.vector_cache:
                # Bat buoc cach ly ngu canh tenant
                if cached_tenant != tenant_id:
                    continue
                sim = self._cosine_similarity(embedding, cached_emb)
                if sim > best_sim:
                    best_sim = sim
                    best_answer = cached_a
                    matched_q = cached_q

            if best_sim >= self.similarity_threshold:
                return best_answer, "SEMANTIC_CACHE_HIT", {
                    "similarity": round(best_sim, 4),
                    "matched_question": matched_q,
                    "latency_ms": 15.0
                }

        # 3. Cache Miss: Can goi LLM thuc te
        return "", "CACHE_MISS", {"hash_key": hash_key}

    def update_cache(
        self, prompt: str, embedding: List[float], response: str, tenant_id: str = "default"
    ):
        """Cap nhat ca hai tang cache khi co phan hoi moi tu LLM."""
        hash_key = self._normalize_and_hash(prompt, tenant_id)
        self.exact_cache[hash_key] = response
        self.vector_cache.append((embedding, prompt, response, tenant_id))

    @staticmethod
    def route_model(prompt: str) -> str:
        """Dinh tuyen phan cap: Chon mo hinh phu hop theo do phuc tap."""
        tokens = prompt.lower().split()
        complex_keywords = {"chung minh", "toan hoc", "kien truc", "debug", "refactor"}
        has_complex_keyword = any(kw in prompt.lower() for kw in complex_keywords)
        is_long = len(tokens) > 100
        
        if has_complex_keyword or is_long:
            return "TIER2_LARGE_MODEL_70B"
        return "TIER1_SMALL_MODEL_8B"
\end{lstlisting}

---

\section{Thiết Kế Hệ Thống: Cân Bằng Tải Nhận Biết Tiền Tố (Prefix-Aware Routing) \& Hạ Tầng Anycast Toàn Cầu}
\label{sec:ch13-prefix-aware-routing}

Trong các kiến trúc microservices truyền thống, thuật toán cân bằng tải (Load Balancer) thường hoạt động theo nguyên lý Round-Robin, Least-Connections, hoặc ngẫu nhiên có trọng số. Tuy nhiên, khi áp dụng các giải thuật này trực tiếp vào cụm phục vụ mô hình ngôn ngữ lớn (vLLM / Triton clusters), toàn bộ hệ thống sẽ rơi vào một cuộc khủng hoảng hiệu năng và lãng phí phần cứng nghiêm trọng.

\subsection{Khởi Điểm Từ Con Số 0: Thảm Họa Prefill Của Cân Bằng Tải Ngây Thơ}
Để hiểu tại sao bộ cân bằng tải cho AI bắt buộc phải nhận biết nội dung tiền tố, ta cần phân tích bản chất hoạt động của bộ đệm \term{PagedAttention} và cấu trúc cây tiền tố \term{RadixAttention} trong vLLM:
\begin{enumerate}
  \item \textbf{Cấu trúc lặp lại của Prompt doanh nghiệp}:
  Trong hầu hết các ứng dụng thực tế (như trợ lý pháp lý, hỗ trợ kỹ thuật khách hàng, hoặc lập trình viên AI), mọi truy vấn đều chứa một phần mở đầu cố định (\textit{System Prompt \& Few-Shot Examples \& Context RAG}) có độ dài rất lớn, thường từ $2.000$ đến $8.000$ tokens. Phần câu hỏi thực tế của người dùng ở cuối câu chỉ chiếm khoảng $50$ tokens.
  \item \textbf{Chi phí tính toán pha Prefill}:
  Khi một máy chủ GPU tiếp nhận một prompt dài $4.000$ tokens mà nó chưa từng nhìn thấy trước đó, GPU bắt buộc phải thực hiện pha \term{Prefill}: Tính toán ma trận Attention toàn phần $\mathcal{O}(L^2)$ trên toàn bộ $4.000$ tokens. Quá trình này tiêu tốn hàng triệu phép tính dấu phẩy động và làm mất từ $300\text{ms}$ đến $500\text{ms}$ thời gian xử lý của GPU.
  \item \textbf{Sức mạnh của Prefix Caching}:
  Nếu máy chủ GPU đó \textbf{đã từng xử lý} đoạn System Prompt $4.000$ tokens này trong các yêu cầu trước, toàn bộ các véc-tơ Key và Value (KV Cache) của $4.000$ tokens đó vẫn đang nằm sẵn trong VRAM của GPU dưới dạng một nhánh cây Radix Tree.
  Khi request mới tới, GPU chỉ cần tái sử dụng lại các block bộ nhớ đã có sẵn: Thời gian Prefill giảm ngoạn mục từ $350\text{ms}$ xuống chỉ còn \textbf{$8\text{ms}$} (giảm hơn $97\%$), và GPU gần như không tiêu tốn thêm một chu kỳ tính toán nào!
\end{enumerate}

\paragraph{Nghịch lý của bộ cân bằng tải ngây thơ (Naive Round-Robin)}:
Nếu bộ cân bằng tải ở Gateway không có tri thức về ngữ cảnh và chia đều các request theo vòng tròn:
\begin{itemize}
  \item Request 1 của Doanh nghiệp X gửi sang GPU Node 1 $\implies$ GPU 1 tốn $350\text{ms}$ tính toán Prefill và lưu cache $4.000$ tokens.
  \item Request 2 của Doanh nghiệp X gửi sang GPU Node 2 $\implies$ GPU 2 lại tốn $350\text{ms}$ tính toán lại từ đầu và nhân bản thêm $4.000$ tokens KV Cache vào VRAM của mình.
  \item Request 3 của Doanh nghiệp X gửi sang GPU Node 3 $\implies$ GPU 3 lại lặp lại chu trình lãng phí tương tự.
\end{itemize}
\textbf{Hậu quả tai hại}: Toàn bộ bộ nhớ VRAM đắt giá của cả cụm GPU bị lấp đầy bởi các bản sao KV Cache trùng lặp giống hệt nhau! Tỷ lệ Cache Hit của toàn cụm tụt dốc về gần $0\%$, thời gian chờ đợi phản hồi TTFT của người dùng tăng vọt gấp hàng chục lần, và doanh nghiệp buộc phải bỏ ra số tiền gấp 5 lần để thuê thêm máy chủ GPU!

\subsection{Mô Hình Toán Học: Consistent Hashing Với Virtual Nodes \& Cân Bằng Knapsack Hai Chiều}
Để giải quyết bài toán này, kiến trúc chuẩn mực của một \term{AI Gateway} hiện đại phải tích hợp thuật toán \term{Prefix-Aware Load Balancer}.

\begin{figure}[htbp]
\centering
\resizebox{0.95\linewidth}{!}{%
\begin{tikzpicture}[node distance=1.2cm, auto,
    client/.style={rectangle, draw=primarycolor, thick, fill=boxbg, align=center, rounded corners=2pt, minimum height=0.7cm, minimum width=2.5cm, font=\scriptsize\bfseries},
    gateway/.style={rectangle, draw=dangercolor, ultra thick, fill=accentcolor!20, align=center, rounded corners=3pt, minimum height=1.0cm, minimum width=11cm, font=\small\bfseries},
    gpu/.style={rectangle, draw=secondarycolor, thick, fill=boxbg, align=center, rounded corners=2pt, minimum height=1.2cm, minimum width=2.4cm, font=\scriptsize\bfseries}
]
    \node[client] (c1) {Client Tenant A};
    \node[client, right=1.5cm of c1] (c2) {Client Tenant A};
    \node[client, right=1.5cm of c2] (c3) {Client Tenant B};
    \node[client, right=1.5cm of c3] (c4) {Client Tenant C};

    \node[gateway, below=0.8cm of c2, xshift=1.2cm] (gw) {AI Gateway: Prefix-Aware Router (Consistent Hashing + Radix Tracker)};

    \node[gpu, below=1.0cm of gw, xshift=-4.0cm] (gpu1) {GPU Node 1\\ \tiny \textbf{Cache: Tenant A}\\ \tiny (Hit Rate: 98\%)};
    \node[gpu, below=1.0cm of gw, xshift=-1.3cm] (gpu2) {GPU Node 2\\ \tiny \textbf{Cache: Tenant A}\\ \tiny (Replica / Failover)};
    \node[gpu, below=1.0cm of gw, xshift=1.4cm] (gpu3) {GPU Node 3\\ \tiny \textbf{Cache: Tenant B}\\ \tiny (Hit Rate: 95\%)};
    \node[gpu, below=1.0cm of gw, xshift=4.1cm] (gpu4) {GPU Node 4\\ \tiny \textbf{Cache: Tenant C}\\ \tiny (Hit Rate: 92\%)};

    \draw[draw=primarycolor, thick, -Stealth] (c1) -- (gw.150);
    \draw[draw=primarycolor, thick, -Stealth] (c2) -- (gw.110);
    \draw[draw=primarycolor, thick, -Stealth] (c3) -- (gw.70);
    \draw[draw=primarycolor, thick, -Stealth] (c4) -- (gw.30);

    \draw[draw=secondarycolor, line width=1.5pt, -Stealth] (gw.210) -- node[left, font=\tiny\bfseries] {Route to Hot Cache} (gpu1);
    \draw[draw=secondarycolor, line width=1.5pt, -Stealth] (gw.330) -- node[right, font=\tiny\bfseries] {Route to Hot Cache} (gpu3);
    \draw[draw=secondarycolor, line width=1.5pt, -Stealth] (gw.350) -- (gpu4);
\end{tikzpicture}%
}
\caption{Kiến trúc Prefix-Aware Load Balancer: Gom các truy vấn cùng tiền tố về đúng máy chủ GPU lưu trữ sẵn KV Cache}
\label{fig:prefix-aware-load-balancer}
\end{figure}

\paragraph{1. Vòng băm Nhất quán (Consistent Hash Ring) cho System Prompt}:
Gateway phân tích nội dung mở đầu của truy vấn và tạo một khóa băm định danh tiền tố:
\begin{equation}
K_{\text{prefix}} = \text{MurmurHash3}\left(\text{tenant\_id} + "\#::\#" + \text{system\_prompt}\right) \pmod{2^{32}}
\end{equation}
Không gian khóa được biểu diễn trên một vòng tròn liên tục từ $0$ đến $2^{32} - 1$.
Mỗi máy chủ GPU $i$ được gán $V = 128$ nút ảo (\textit{Virtual Nodes}) phân bổ đều trên vòng băm để đảm bảo tính phân tán ngẫu nhiên đồng đều:
\begin{equation}
\text{vnode}_{i, j} = \text{MD5}\left(\text{"Server\_"} + i + "\#vn\_" + j\right) \pmod{2^{32}}
\end{equation}
Khi một truy vấn đi vào, Gateway tìm kiếm nút ảo đầu tiên có vị trí lớn hơn hoặc bằng $K_{\text{prefix}}$ theo chiều kim đồng hồ.
\begin{theorem}[Bảo Toàn Bộ Đệm Khi Mở Rộng Cụm]
Khi thêm một máy chủ GPU mới vào cụm gồm $N$ máy chủ hoặc khi một máy chủ gặp sự cố, tỷ lệ khóa tiền tố bị dịch chuyển sang máy chủ khác chỉ bằng đúng:
\begin{equation}
\Delta_{\text{keys}} = \frac{1}{N}
\end{equation}
Toàn bộ $\frac{N-1}{N}$ ($> 90\%$) các khóa tiền tố còn lại vẫn giữ nguyên vị trí, bảo toàn tuyệt đối tỷ lệ Cache Hit của toàn bộ hệ thống!
\end{theorem}

\paragraph{2. Giải thuật điều hướng cân bằng hai chiều (Dual-Constraint Routing Score)}:
Nếu chỉ dựa hoàn toàn vào Consistent Hashing, hệ thống sẽ đối mặt với rủi ro \term{Hotspot Node}: Một khách hàng doanh nghiệp lớn với hàng ngàn nhân viên gửi truy vấn dồn dập có thể làm nghẽn hoàn toàn máy chủ GPU phụ trách tiền tố đó, trong khi các máy chủ khác vẫn rảnh rỗi.

Để giải quyết, Gateway áp dụng hàm chấm điểm đa mục tiêu hai chiều:
\begin{equation}
\text{Score}(i) = \alpha \cdot \left(\frac{\text{SharedPrefixTokens}(P, \mathcal{T}_i)}{L_{\text{prompt}}}\right) - (1 - \alpha) \cdot \left(\frac{Q_i}{C_i}\right)
\end{equation}
Trong đó:
\begin{itemize}
  \item $\text{SharedPrefixTokens}(P, \mathcal{T}_i)$: Số lượng tokens tiền tố đang có sẵn trong VRAM của GPU $i$.
  \item $Q_i$: Số lượng request đang nằm trong hàng đợi chờ xử lý của GPU $i$.
  \item $C_i$: Ngưỡng dung lượng hàng đợi tối đa cho phép của GPU $i$.
  \item $\alpha \in [0, 1]$: Trọng số đánh đổi (thực nghiệm chọn $\alpha = 0.75$).
\end{itemize}
Khi hàng đợi của GPU nắm giữ cache vượt quá ngưỡng an toàn ($Q_i \ge 0.85 C_i$), số hạng phạt thứ hai sẽ áp đảo số hạng thứ nhất, buộc Gateway phải chuyển hướng yêu cầu sang máy chủ GPU dự phòng ít tải hơn, bảo đảm cân bằng giữa tối đa hóa Cache Hit và triệt tiêu hàng đợi tắc nghẽn.

\subsection{Kiến Trúc Mạng Anycast BGP Toàn Cầu \& Cụm Suy Luận Đa Vùng}
Để phục vụ người dùng trên toàn cầu với độ trễ mạng thấp nhất, hệ thống kết hợp kiến trúc định tuyến \term{BGP Anycast}:
\begin{enumerate}
  \item \textbf{Một địa chỉ IP duy nhất cho toàn thế giới}: Tất cả các trung tâm dữ liệu tại Bắc Mỹ (us-east), Châu Âu (eu-west) và Châu Á (ap-southeast) đều cùng công bố một dải địa chỉ IP công khai duy nhất (\textit{Anycast IP}) tới các nhà mạng viễn thông toàn cầu qua giao thức BGP.
  \item \textbf{Định tuyến chặng ngắn nhất (Shortest AS-Path)}: Bộ định tuyến Internet tự động dẫn luồng gói tin của người dùng tại Việt Nam vào cụm máy chủ tại Singapore (độ trễ mạng $< 30\text{ms}$), thay vì bay nửa vòng trái đất sang Mỹ ($250\text{ms}$).
  \item \textbf{Tự động khôi phục sự cố tuyến cáp}: Nếu trung tâm dữ liệu tại Singapore gặp sự cố mất điện, bộ định tuyến BGP tự động thu hồi tuyến Anycast trong vài giây, các gói tin của người dùng lập tức được chuyển hướng thông suốt sang trung tâm dữ liệu tại Nhật Bản hoặc Châu Âu mà người dùng không hề nhận thấy sự gián đoạn.
\end{enumerate}

---

\section{Kiến Trúc Định Tuyến GPU Đa Vùng Toàn Cầu \& Khôi Phục Thảm Họa (Multi-Region Geo-Distributed Routing \& Disaster Recovery)}
\label{sec:ch13-multi-region-routing}

Đối với các ứng dụng AI quy mô toàn cầu phục vụ hàng chục triệu người dùng (như Canva, Stripe, hay Grab), việc triển khai cụm suy luận tại một trung tâm dữ liệu duy nhất (\textit{Single-Region Deployment}) tạo ra hai rủi ro sống còn:
\begin{itemize}
  \item \textbf{Độ trễ truyền dẫn xuyên lục địa (Speed of Light Latency)}: Gói tin truyền trong cáp quang biển giữa Đông Nam Á và Bắc Mỹ mất tối thiểu $220 - 250\text{ms}$ cho một vòng khứ hồi (RTT). Với các ứng dụng Agent tương tác nhiều bước (Multi-turn Tool Calling), độ trễ mạng này sẽ cộng dồn thành nhiều giây chờ đợi khó chịu.
  \item \textbf{Điểm nghẽn lỗi đơn độc (Single Point of Failure -- SPoF)}: Mất điện diện rộng, đứt cáp quang biển, hoặc trung tâm dữ liệu của nhà cung cấp điện toán đám mây gặp sự cố có thể làm tê liệt $100\%$ dịch vụ của doanh nghiệp.
\end{itemize}

\subsection{Mô Hình Định Tuyến Trọng Số Trễ Động (Dynamic Latency-Weighted Routing)}
Hạ tầng AI toàn cầu thường được bố trí tại 3 cụm khu vực trọng điểm: \textbf{AP-Southeast} (Singapore), \textbf{EU-West} (Frankfurt), và \textbf{US-East} (Virginia).
Độ trễ mạng cơ sở (\textit{Base Network RTT}) giữa người dùng tại Việt Nam và các cụm này lần lượt là:
\begin{equation}
L_{\text{base}}(\text{Singapore}) \approx 25\text{ms}, \quad L_{\text{base}}(\text{Frankfurt}) \approx 160\text{ms}, \quad L_{\text{base}}(\text{Virginia}) \approx 220\text{ms}
\end{equation}

Tuy nhiên, nếu chỉ định tuyến thuần túy theo độ trễ mạng thấp nhất, cụm máy chủ Singapore sẽ nhanh chóng bị bão hòa tải vào khung giờ cao điểm của Châu Á, dẫn đến hiện tượng hàng đợi GPU phình to.
Để tối ưu hóa toàn diện, Gateway tính toán \textbf{Độ Trễ Hiệu Dụng (Effective Latency)}:
\begin{equation}
L_{\text{eff}}(r) = L_{\text{base}}(r) \times \left(1 + \frac{\text{ActiveLoad}(r)}{\text{Capacity}(r)}\right)
\label{eq:effective-latency}
\end{equation}
Trong đó $\text{ActiveLoad}(r)$ là số yêu cầu đang nằm trong hàng đợi xử lý của vùng $r$, và $\text{Capacity}(r)$ là năng lực phục vụ tối đa của cụm GPU vùng đó.
Khi cụm Singapore bị quá tải $100\%$ ($\text{ActiveLoad} \approx \text{Capacity}$), $L_{\text{eff}}(\text{Singapore})$ tăng vọt từ $25\text{ms}$ lên $50\text{ms}$ (hoặc hàng trăm ms nếu nghẽn hàng đợi sâu). Khi đó, Gateway có thể tự động san sẻ tải sang các cụm lân cận có độ trễ cơ sở lớn hơn nhưng VRAM đang rảnh rỗi hoàn toàn.

\subsection{Máy Trạng Thái Ngắt Mạch \& Thác Phục Hồi Thảm Họa (Circuit Breaker \& Graceful Fallback)}
Để đảm bảo cam kết tính sẵn sàng $99.99\%$, mỗi kết nối tới một vùng GPU được giám sát bởi một \term{Bộ Ngắt Mạch (Circuit Breaker)} với máy trạng thái 3 pha:
\begin{figure}[htbp]
\centering
\begin{tikzpicture}[node distance=2.2cm, auto,
    state/.style={circle, draw=primarycolor, thick, fill=boxbg, align=center, minimum size=1.8cm, font=\small\bfseries},
    arrow/.style={draw=primarycolor, -Stealth, thick}
]
    \node[state, fill=green!15] (closed) {CLOSED\\ \tiny (Normal)};
    \node[state, right=2.5cm of closed, fill=red!15] (open) {OPEN\\ \tiny (Tripped)};
    \node[state, below=1.8cm of $(closed)!0.5!(open)$, fill=yellow!20] (half) {HALF-OPEN\\ \tiny (Testing)};

    \path[arrow] (closed) edge[bend left=20] node[above, font=\scriptsize] {$N_{\text{fail}} \ge 3$ consecutive errors} (open);
    \path[arrow] (open) edge[bend left=20] node[below right, font=\scriptsize] {$T_{\text{cooldown}} \ge 30\text{s}$} (half);
    \path[arrow] (half) edge[bend left=20] node[below left, font=\scriptsize] {$N_{\text{success}} \ge 2$ test passed} (closed);
    \path[arrow] (half) edge[bend right=20] node[right, font=\scriptsize] {Any test error} (open);
\end{tikzpicture}
\caption{Máy trạng thái 3 pha Circuit Breaker bảo vệ cụm suy luận GPU phân tán}
\label{fig:circuit-breaker-state-machine}
\end{figure}

\begin{enumerate}
  \item \textbf{Trạng thái CLOSED (Bình thường)}: Mọi truy vấn được gửi tới cụm GPU mục tiêu. Nếu xảy ra $N_{\text{fail}} \ge 3$ lần lỗi liên tiếp (mã lỗi HTTP 500, 503 hoặc Timeout quá $3.0\text{ giây}$), mạch ngắt lập tức nhảy sang trạng thái \textbf{OPEN}.
  \item \textbf{Trạng thái OPEN (Ngắt mạch hoàn toàn)}: Mọi yêu cầu tới vùng này bị chặn đứng ngay lập tức tại Gateway mà không cần gửi gói tin thử nghiệm qua mạng. Luồng dữ liệu được tự động kích hoạt thác phục hồi (\textit{Graceful Fallback}):
  \begin{align*}
  \text{Primary On-Premise GPU} &\longrightarrow \text{Secondary Geo-Region GPU} \\
  &\longrightarrow \text{Cloud Provider Fallback (Bedrock/Azure)}
  \end{align*}
  \item \textbf{Trạng thái HALF-OPEN (Kiểm định phục hồi)}: Sau khoảng thời gian chờ nguội $T_{\text{cooldown}} = 30\text{ giây}$, mạch ngắt chuyển sang trạng thái HALF-OPEN và cho phép một tỷ lệ lưu lượng nhỏ ($5\%$) đi qua để thăm dò. Nếu cụm máy chủ đáp ứng thành công $N_{\text{success}} \ge 2$ lần liên tiếp, mạch đóng lại (\textbf{CLOSED}); nếu xuất hiện bất kỳ lỗi nào, mạch nhảy ngược lại về \textbf{OPEN}.
\end{enumerate}

---

\section{Quản Trị Lưu Lượng Phân Tán, Rate Limiting \& Ngân Sách Sai Số SLO/SLA}
\label{sec:ch13-rate-limiting-slo}

Trong môi trường doanh nghiệp đa khách hàng (\textit{Multi-Tenant SaaS}), một khách hàng đơn lẻ gửi truy vấn quá mức có thể làm cạn kiệt toàn bộ tài nguyên VRAM của cụm GPU, gây ảnh hưởng tiêu cực tới hàng ngàn khách hàng khác (\textit{Noisy Neighbor Problem}). Quản trị lưu lượng phân tán là tuyến phòng thủ bắt buộc.

\subsection{Giải Thuật Xô Thẻ Phân Tán (Distributed Token Bucket via Redis Lua)}
So với thuật toán Cửa Sổ Cố Định (Fixed Window Counter vốn dễ bị lỗi bùng nổ lưu lượng tại biên giới hạn) hay Cửa Sổ Trượt (Sliding Window Log tiêu tốn nhiều bộ nhớ RAM), thuật toán \term{Token Bucket} là chuẩn mực công nghiệp vàng vì hỗ trợ tính năng bùng nổ lưu lượng tức thời (\textit{Burst Capacity}) mà vẫn kiểm soát chặt chẽ tốc độ dòng chảy trung bình (\textit{Sustained Rate}).

\paragraph{Công thức toán học của Token Bucket}:
Xét một khách hàng doanh nghiệp được cấp hạn mức:
\begin{itemize}
  \item Tốc độ nạp thẻ: $R$ (tokens/giây, tương đương số yêu cầu tối đa mỗi giây ở chế độ bền vững).
  \item Dung lượng tối đa của xô: $B$ (Burst Limit, số yêu cầu tối đa được phép thực hiện tức thời).
\end{itemize}
Tại thời điểm yêu cầu mới tới $t$, số lượng thẻ hiện có trong xô được nạp bù theo khoảng thời gian đã trôi qua kể từ lần nạp trước $t_{\text{last}}$:
\begin{equation}
T(t) = \min\Big( B, \; T(t_{\text{last}}) + R \cdot (t - t_{\text{last}}) \Big)
\label{eq:token-bucket-refill}
\end{equation}
Nếu số lượng thẻ yêu cầu $k$ thỏa mãn $k \le T(t)$: Yêu cầu được chấp thuận, và số thẻ mới là $T(t) - k$.
Nếu $k > T(t)$: Yêu cầu bị từ chối với mã lỗi \code{HTTP 429 Too Many Requests}, kèm theo tiêu đề phản hồi thời gian chờ tối thiểu để thử lại (\textit{Retry-After}):
\begin{equation}
\Delta t_{\text{retry}} = \frac{k - T(t)}{R}
\label{eq:token-bucket-retry}
\end{equation}
Trong môi trường phân tán gồm hàng chục Gateway instances, thao tác kiểm tra và cập nhật thẻ bắt buộc phải được đóng gói trong một script \textbf{Redis Lua} duy nhất để đảm bảo tính nguyên tử tuyệt đối (Atomicity $\mathcal{O}(1)$), tránh xung đột dữ liệu (Race Condition).

\subsection{Phân Tích Hàng Đợi M/M/c \& Thảm Họa Độ Trễ Đuôi (Tail Latency Explosion)}
Mô hình toán học kinh điển phân tích hiệu năng phục vụ của cụm máy chủ GPU là lý thuyết hàng đợi \term{M/M/c} (Poisson Arrival, Exponential Service Time, $c$ servers):
\begin{itemize}
  \item Tốc độ luồng yêu cầu gửi tới cụm: $\lambda$ (requests/giây).
  \item Tốc độ xử lý của một nút GPU: $\mu$ (requests/giây).
  \item Số nút GPU phục vụ độc lập trong cụm: $c$.
  \item Hệ số cường độ tải của cụm:
  \begin{equation}
  \rho = \frac{\lambda}{c \cdot \mu} < 1
  \end{equation}
\end{itemize}
Theo công thức Erlang-C, thời gian chờ trung bình trong hàng đợi $W_q$ xấp xỉ:
\begin{equation}
W_q \approx \frac{C(c, \lambda/\mu)}{c\mu (1 - \rho)} \propto \frac{1}{1 - \rho}
\end{equation}
\textbf{Hiện tượng bùng nổ độ trễ đuôi (Tail Latency Degeneration)}:
Khi hệ số tải $\rho \le 0.70$ ($70\%$ công suất), hàng đợi hầu như rỗng, độ trễ $P50$ và $P99$ bám sát thời gian tính toán thực tế của GPU ($\approx 300\text{ms}$).
Tuy nhiên, khi tải vượt ngưỡng $85\% \to 95\%$ ($\rho \to 1$), mẫu số $1 - \rho \to 0$ khiến thời gian chờ trong hàng đợi $W_q$ tăng phi tuyến theo đường tiệm cận đứng!
Tại phân vị $P99$ và $P99.9$, thời gian phản hồi tăng vọt từ $300\text{ms}$ lên tới $15 - 30\text{ giây}$, phá vỡ hoàn toàn cam kết chất lượng dịch vụ của doanh nghiệp. Do đó, quy tắc thiết kế hệ thống tối thượng là: \textbf{Không bao giờ để cụm GPU vận hành liên tục vượt quá $75 - 80\%$ tải thiết kế mà không kích hoạt chính sách cắt tải (Load Shedding)!}

\subsection{Ngân Sách Sai Số SLO/SLA \& Tốc Độ Đốt (Error Budget \& Burn Rate)}
Để cân bằng giữa tốc độ triển khai tính năng mới và tính ổn định hệ thống, chuẩn mực kỹ nghệ SRE của Google đưa ra khái niệm \term{Ngân Sách Sai Số (Error Budget)}:
\begin{equation}
\text{Error Budget} = 1.0 - \text{SLO Target}
\end{equation}
Với cam kết chất lượng dịch vụ $\text{SLA} = 99.9\%$ ($3$ số 9) trong chu kỳ 30 ngày ($2.592.000$ giây):
\begin{itemize}
  \item Tỷ lệ lỗi tối đa được phép: $1 - 0.999 = 0.001$ ($0.1\%$).
  \item Thời gian chết tối đa cho phép (\textit{Downtime Budget}):
  \begin{equation}
  T_{\text{downtime}} = 2.592.000 \times 0.001 = 2.592 \text{ giây} \approx \mathbf{43.2 \text{ phút/tháng}}
  \end{equation}
  \item Với một hệ thống tiếp nhận $100.000.000$ truy vấn mỗi tháng, số truy vấn lỗi tối đa được phép chấp nhận là đúng $100.000$ requests.
\end{itemize}

\paragraph{Chỉ số Tốc Độ Đốt Ngân Sách (Burn Rate)}:
Burn Rate đo lường tốc độ tiêu thụ ngân sách lỗi so với tốc độ tiêu thụ dự kiến trong cả tháng:
\begin{equation}
\text{Burn Rate} = \frac{\text{Observed Error Rate}}{1 - \text{SLA Target}} = \frac{\text{Observed Error Rate}}{0.001}
\label{eq:burn-rate-formula}
\end{equation}
Quy tắc kích hoạt cảnh báo khẩn cấp theo chuẩn Google SRE:
\begin{itemize}
  \item $\text{Burn Rate} = 1.0$: Hệ thống đang tiêu hao vừa đúng $100\%$ ngân sách trong 30 ngày. An toàn.
  \item $\text{Burn Rate} = 6.0$: Hệ thống tiêu hao hết $5\%$ ngân sách cả tháng chỉ trong vòng $6$ giờ $\implies$ Tự động tạo Ticket cảnh báo mức cao cho đội ngũ vận hành.
  \item $\text{Burn Rate} = 14.4$: Hệ thống tiêu hao hết $2\%$ ngân sách cả tháng chỉ trong vòng $1$ giờ đồng hồ (tương ứng tỷ lệ lỗi quan sát là $1.44\%$) $\implies$ Kích hoạt chuông báo động khẩn cấp \textbf{Critical Pager Alert}, đánh thức kỹ sư trực hệ thống On-Call ngay lập tức!
\end{itemize}

---

\section{Hiện Thực Hóa: Mã Nguồn Cài Đặt Chuẩn Kỹ Nghệ \& Prompts Hoạt Họa 60 FPS}
\label{sec:ch13-advanced-code}

Dưới đây là các lớp mã nguồn Python kỹ nghệ hoàn chỉnh, không phụ thuộc thư viện ngoài, được kiểm thử $100\%$ và có thể tích hợp trực tiếp vào hạ tầng AI Gateway thực tế:

\subsection{1. Cài Đặt Bộ Kiểm Soát Lưu Lượng Token Bucket Phân Tán}
\begin{lstlisting}[language=Python]
from typing import Dict, Any, Tuple, Optional

class DistributedTokenBucketRateLimiter:
    """
    Distributed Token Bucket Rate Limiter supporting burst capacity,
    multi-tenant tiers, and atomic token refill logic (simulating Redis Lua).
    """
    def __init__(self, default_rate: float = 10.0, default_capacity: float = 20.0):
        self.default_rate = default_rate
        self.default_capacity = default_capacity
        self.tenants: Dict[str, Dict[str, Any]] = {}

    def register_tenant(self, tenant_id: str, rate: float, capacity: float):
        self.tenants[tenant_id] = {
            "rate": float(rate),
            "capacity": float(capacity),
            "tokens": float(capacity),
            "last_refill": None,
        }

    def _get_tenant(self, tenant_id: str) -> Dict[str, Any]:
        if tenant_id not in self.tenants:
            self.register_tenant(tenant_id, self.default_rate, self.default_capacity)
        return self.tenants[tenant_id]

    def acquire(
        self,
        tenant_id: str,
        tokens_requested: int = 1,
        current_time: float = 0.0
    ) -> Tuple[bool, float, Dict[str, Any]]:
        tenant = self._get_tenant(tenant_id)
        rate = tenant["rate"]
        capacity = tenant["capacity"]
        last_refill = tenant["last_refill"]
        tokens = tenant["tokens"]

        # 1. Refill tokens based on elapsed time
        if last_refill is not None and current_time > last_refill:
            elapsed = current_time - last_refill
            tokens = min(capacity, tokens + elapsed * rate)

        tenant["last_refill"] = current_time

        # 2. Check if sufficient tokens are available
        if tokens >= tokens_requested:
            tokens -= tokens_requested
            tenant["tokens"] = tokens
            return True, 0.0, {
                "tenant_id": tenant_id,
                "allowed": True,
                "remaining_tokens": round(tokens, 2),
                "retry_after_sec": 0.0,
            }
        else:
            needed = tokens_requested - tokens
            retry_after = needed / rate if rate > 0 else float("inf")
            tenant["tokens"] = tokens
            return False, round(retry_after, 3), {
                "tenant_id": tenant_id,
                "allowed": False,
                "remaining_tokens": round(tokens, 2),
                "retry_after_sec": round(retry_after, 3),
            }
\end{lstlisting}

\subsection{2. Cài Đặt Bộ Định Tuyến GPU Đa Vùng Kèm Ngắt Mạch}
\begin{lstlisting}[language=Python]
class MultiRegionGPURouter:
    """
    Multi-Region Geo-Distributed GPU Router with dynamic latency weighting,
    load-aware penalty, circuit breaker (Closed/Open/Half-Open), and graceful cloud fallback.
    """
    def __init__(self, failure_threshold: int = 3, recovery_threshold: int = 2):
        self.failure_threshold = failure_threshold
        self.recovery_threshold = recovery_threshold
        self.regions: Dict[str, Dict[str, Any]] = {}

    def add_region(self, name: str, base_latency_ms: float, capacity: int):
        self.regions[name] = {
            "name": name,
            "base_latency_ms": base_latency_ms,
            "capacity": capacity,
            "active_load": 0,
            "consecutive_failures": 0,
            "consecutive_successes": 0,
            "circuit_state": "CLOSED",  # CLOSED, OPEN, HALF_OPEN
        }

    def record_result(self, region_name: str, success: bool):
        if region_name not in self.regions:
            return
        reg = self.regions[region_name]
        if success:
            reg["consecutive_failures"] = 0
            reg["consecutive_successes"] += 1
            if reg["circuit_state"] == "HALF_OPEN" and reg["consecutive_successes"] >= self.recovery_threshold:
                reg["circuit_state"] = "CLOSED"
        else:
            reg["consecutive_successes"] = 0
            reg["consecutive_failures"] += 1
            if reg["consecutive_failures"] >= self.failure_threshold:
                reg["circuit_state"] = "OPEN"

    def route(
        self,
        tenant_id: str = "default",
        max_acceptable_latency_ms: float = 300.0
    ) -> Tuple[str, str, Dict[str, Any]]:
        candidates = []
        for reg in self.regions.values():
            if reg["circuit_state"] in ("CLOSED", "HALF_OPEN"):
                load_factor = reg["active_load"] / max(1, reg["capacity"])
                effective_latency = reg["base_latency_ms"] * (1.0 + load_factor)
                candidates.append((effective_latency, reg))

        if not candidates:
            return "CLOUD_PROVIDER_FALLBACK", "CIRCUIT_BREAKER_ALL_OPEN", {
                "reason": "All on-premise GPU clusters unreachable or tripped",
                "effective_latency_ms": 250.0,
            }

        candidates.sort(key=lambda x: x[0])
        best_eff_latency, best_reg = candidates[0]

        if best_eff_latency > max_acceptable_latency_ms:
            return "CLOUD_PROVIDER_FALLBACK", "LATENCY_SLA_VIOLATION", {
                "reason": f"Lowest internal latency {best_eff_latency:.1f}ms exceeds SLA {max_acceptable_latency_ms}ms",
                "effective_latency_ms": best_eff_latency,
            }

        best_reg["active_load"] += 1
        return best_reg["name"], "ROUTED_OPTIMAL_CLUSTER", {
            "effective_latency_ms": round(best_eff_latency, 2),
            "circuit_state": best_reg["circuit_state"],
            "active_load": best_reg["active_load"],
        }
\end{lstlisting}

\subsection{3. Cài Đặt Bộ Đo Ngân Sách Lỗi SLO \& Phân Vị Độ Trễ}
\begin{lstlisting}[language=Python]
import math
from typing import List

class SLOErrorBudgetCalculator:
    """
    Mathematical modeling for SLO/SLA error budgets, burn rate alerts,
    and tail latency percentiles (P50, P90, P99, P99.9).
    """
    def __init__(self, sla_target: float = 0.999):
        self.sla_target = sla_target
        self.allowed_error_fraction = 1.0 - sla_target

    def calculate_budget(self, total_requests: int, failed_requests: int) -> Dict[str, Any]:
        allowed_failures = math.floor(total_requests * self.allowed_error_fraction)
        remaining_failures = allowed_failures - failed_requests
        budget_used_pct = (failed_requests / allowed_failures * 100.0) if allowed_failures > 0 else 0.0
        return {
            "total_requests": total_requests,
            "failed_requests": failed_requests,
            "allowed_failures": allowed_failures,
            "remaining_allowed_failures": max(0, remaining_failures),
            "budget_used_percentage": round(budget_used_pct, 2),
            "budget_exhausted": failed_requests > allowed_failures,
        }

    def compute_burn_rate(self, observed_error_rate: float) -> float:
        if self.allowed_error_fraction <= 0.0:
            return float("inf")
        return round(observed_error_rate / self.allowed_error_fraction, 2)

    @staticmethod
    def calculate_percentiles(latencies_ms: List[float]) -> Dict[str, float]:
        if not latencies_ms:
            return {"p50": 0.0, "p90": 0.0, "p99": 0.0, "p99_9": 0.0}
        sorted_vals = sorted(latencies_ms)
        n = len(sorted_vals)

        def get_p(p: float) -> float:
            k = (n - 1) * p
            f = math.floor(k)
            c = math.ceil(k)
            if f == c:
                return sorted_vals[int(k)]
            d0 = sorted_vals[int(f)] * (c - k)
            d1 = sorted_vals[int(c)] * (k - f)
            return d0 + d1

        return {
            "p50": round(get_p(0.50), 2),
            "p90": round(get_p(0.90), 2),
            "p99": round(get_p(0.99), 2),
            "p99_9": round(get_p(0.999), 2),
        }
\end{lstlisting}

\begin{simulationbox}[Prompts AI Mô Phỏng Hoạt Họa 60 FPS: Manim CE Trực Quan Hóa Định Tuyến Đa Vùng \& Rate Limiting]
Dưới đây là 2 kịch bản prompt chi tiết cho các công cụ AI tạo mã hoạt họa Manim CE để trực quan hóa sinh động cơ chế định tuyến GPU đa vùng và kiểm soát tốc độ lưu lượng:

\paragraph{Prompt 1: Trực Quan Hóa Định Tuyến GPU Đa Vùng Toàn Cầu \& Ngắt Mạch Tự Động}:
\begin{quote}
\small\itshape
``Hãy đóng vai một chuyên gia Manim CE (Community Edition) và Kiến trúc sư Hạ tầng AI Toàn cầu. \\
Hãy viết một script Python hoàn chỉnh sử dụng thư viện \code{manim} để tạo video hoạt họa 1080p 60fps trực quan hóa mạng lưới định tuyến GPU đa vùng:
\begin{enumerate}
  \item Vẽ quả địa cầu phẳng với 3 cụm GPU: Singapore (màu xanh lá), Frankfurt (màu xanh dương), và Virginia (màu tím). Hiển thị thanh đo dung lượng hàng đợi cho từng cụm.
  \item Một luồng các khối hộp request phát ra từ người dùng tại Việt Nam đi vào khối AI Gateway ở giữa.
  \item Giai đoạn 1 (Bình thường): Gateway tính toán $L_{\text{eff}}$ và định tuyến luồng dữ liệu về cụm Singapore (độ trễ hiển thị: $25\text{ms}$).
  \item Giai đoạn 2 (Quá tải \& Ngắt mạch): Cụm Singapore nhận tải dồn dập, hàng đợi đầy tràn màu đỏ ($100\%$). Xuất hiện 3 tia sét biểu thị lỗi kết nối liên tiếp. Biểu tượng Circuit Breaker nhảy từ trạng thái xanh CLOSED sang đỏ rực OPEN.
  \item Giai đoạn 3 (Tự động chuyển hướng): Cổng vào Singapore đóng sập lại. Gateway chuyển hướng luồng dữ liệu sang cụm Frankfurt (độ trễ $160\text{ms}$). Các khối hộp tiếp tục được xử lý trơn tru mà không có request nào bị mất.
  \item Render bằng lệnh: \code{manim -pql scene.py MultiRegionGPURoutingScene}.''
\end{enumerate}
\end{quote}

\paragraph{Prompt 2: Trực Quan Hóa Thuật Toán Token Bucket Rate Limiting \& Bùng Nổ Độ Trễ Đuôi}:
\begin{quote}
\small\itshape
``Hãy đóng vai một chuyên gia Manim CE và Kỹ sư SRE Cấp cao. \\
Hãy viết một script Python hoàn chỉnh sử dụng thư viện \code{manim} để tạo video hoạt họa 1080p 60fps mô phỏng thuật toán Token Bucket và sự suy thoái độ trễ đuôi $P99$:
\begin{enumerate}
  \item Vẽ một chiếc xô trong suốt (Bucket) có vạch chia dung lượng tối đa $B = 20$ thẻ màu vàng óng. Phía trên có vòi nước nhỏ giọt đều đặn nạp thẻ với tốc độ $R = 10\text{ thẻ/giây}$.
  \item Bên dưới là cổng kiểm soát Gateway. Người dùng gửi đến một đợt sóng bùng nổ 15 requests đồng thời. 15 thẻ rơi qua cổng thành công, trong xô còn lại 5 thẻ.
  \item Đợt sóng thứ hai gồm 10 requests ập tới ngay lập tức: 5 requests đầu lấy hết 5 thẻ còn lại và được thông qua; 5 requests sau chạm đáy xô rỗng, bị cổng chắn màu đỏ đẩy bật ra kèm nhãn \code{HTTP 429: Too Many Requests (Retry-After: 0.5s)}.
  \item Nửa bên phải màn hình: Đồ thị đường cong Erlang-C thể hiện mối quan hệ giữa hệ số tải $\rho$ và thời gian chờ đợi $W_q$. Khi $\rho$ tăng từ $0.70$ lên $0.95$, thanh đo độ trễ $P99$ vọt lên gấp 20 lần kèm theo tiếng chuông báo động PagerDuty.
  \item Render bằng lệnh: \code{manim -pql scene.py TokenBucketRateLimiterScene}.''
\end{enumerate}
\end{quote}
\end{simulationbox}

---

\section{Hệ Thống Bài Tập Kỹ Nghệ \& Đánh Giá Năng Lực}
\label{sec:ch13-exercises}

\subsection{Nhóm A: Câu Hỏi Tư Duy Bản Chất (Deep Intuition)}

\begin{enumerate}[label=\textbf{A\arabic*.}]
  \item \textbf{Phân tích rủi ro rò rỉ dữ liệu qua Semantic Cache}: Tại sao việc triển khai bộ nhớ đệm ngữ nghĩa (Semantic Caching) có thể gây ra rủi ro vi phạm quyền riêng tư và phân tách dữ liệu đa khách hàng (Multi-Tenant Data Isolation) trong hệ thống SaaS? Nêu giải pháp kiến trúc để khắc phục triệt để.
  \item \textbf{Đánh đổi kích thước mô hình Embedding}: Phân tích sự đánh đổi khi lựa chọn kích thước mô hình embedding cho Semantic Cache: Nếu ta dùng một mô hình embedding lớn ($1.536$ chiều, độ trễ $40\text{ms}$) so với một mô hình nhỏ ($384$ chiều, độ trễ $4\text{ms}$), tổng độ trễ của hệ thống trong trường hợp Cache Miss bị ảnh hưởng ra sao?
  \item \textbf{Chỉ số đánh giá trải nghiệm người dùng}: Tại sao trong cụm máy chủ LLM phục vụ đa người dùng, ta lại ưu tiên chỉ số TPOT (Time Per Output Token) hơn là chỉ số tổng thời gian phản hồi (Total Latency) khi đánh giá trải nghiệm người dùng cuối?
  \item \textbf{Nguy cơ Cascading Storm}: Trong kiến trúc định tuyến phân cấp (SLM kết hợp LLM), hiện tượng gì sẽ xảy ra nếu mô hình nhỏ 8B đồng loạt gặp dữ liệu lạ (\textit{Out-of-Distribution}) khiến tỷ lệ Fallback lên mô hình 70B tăng đột biến từ $10\%$ lên $90\%$? Giải pháp bảo vệ cụm máy chủ là gì?
  \item \textbf{Bảo toàn Cache của Consistent Hashing}: Tại sao thuật toán Consistent Hashing với Virtual Nodes lại bảo toàn tỷ lệ Cache Hit của cụm vLLM khi xảy ra sự cố sập một GPU node, trong khi hàm băm chia dư ngây thơ $H(k) \pmod N$ lại làm xóa sạch $100\%$ bộ nhớ đệm KV Cache?
  \item \textbf{Ổn định kết nối Anycast BGP}: Trong kiến trúc định tuyến toàn cầu Anycast BGP cho các cụm AI phục vụ suy luận, hiện tượng Route Flapping ảnh hưởng như thế nào đến kết nối HTTP/2 Streaming SSE của người dùng? Nêu giải pháp khắc phục ở tầng Transport.
  \item \textbf{Cơ Chế Bùng Nổ Độ Trễ Đuôi ($P99/P99.9$) Dưới Góc Nhìn Lý Thuyết Hàng Đợi M/M/c}: Phân tích tại sao trong một cụm GPU phục vụ LLM, khi hệ số sử dụng tài nguyên $\rho$ tăng từ $0.70$ lên $0.95$, thời gian chờ trong hàng đợi $W_q$ không tăng tuyến tính mà lại bùng nổ theo đường tiệm cận đứng $\frac{1}{1-\rho}$? Hiện tượng này làm trầm trọng hóa chỉ số $P99$ và $P99.9$ ra sao? Tại sao một kiến trúc sư hệ thống không bao giờ được cấu hình hệ thống tự động co giãn (Autoscaling) kích hoạt ở ngưỡng $90\%$ GPU utilization?
  \item \textbf{Chiến Lược Ngắt Mạch Circuit Breaker \& Fallback Đa Tầng Cho Cụm GPU Toàn Cầu}: Phân tích hoạt động của 3 trạng thái trong mô hình Circuit Breaker (CLOSED, OPEN, HALF-OPEN) khi áp dụng cho kết nối giữa AI Gateway và cụm GPU đa vùng. Tại sao trạng thái HALF-OPEN lại bắt buộc phải có thời gian chờ nguội $T_{\text{cooldown}}$ và chỉ cho phép một tỷ lệ nhỏ lưu lượng ($5\%$) đi qua thăm dò thay vì mở lại $100\%$ ngay lập tức?
  \item \textbf{Ý Nghĩa Toán Học Của Tốc Độ Đốt Ngân Sách (Burn Rate) Trong Quản Trị SLO/SLA}: Giải thích tại sao trong kỹ nghệ SRE hiện đại của Google, người ta không đợi đến khi hết chu kỳ 30 ngày mới đánh giá vi phạm SLA, mà phải theo dõi liên tục chỉ số Burn Rate theo nhiều cửa sổ thời gian (1 giờ và 6 giờ)? Tại sao mức $\text{Burn Rate} = 14.4\times$ lại được chọn làm ngưỡng kích hoạt chuông báo động khẩn cấp (Critical Pager Alert)?
\end{enumerate}

\subsection{Nhóm B: Tính Toán Giải Tích \& Số Học Bằng Tay}

\begin{enumerate}[label=\textbf{B\arabic*.}]
  \item \textbf{Tính toán ROI \& Chi phí tiết kiệm khi có Semantic Caching}:
  Một công ty công nghệ vận hành hệ thống AI tiếp nhận $5.000.000$ truy vấn mỗi ngày.
  Chi phí trung bình khi gọi mô hình LLM chính là $\$0.01$ cho mỗi truy vấn.
  Hệ thống triển khai bộ nhớ đệm hai tầng với các thông số đo đạc thực tế:
  \begin{itemize}
    \item Tỷ lệ Exact Hash Cache Hit: $15\%$
    \item Tỷ lệ Semantic Vector Cache Hit: $25\%$
    \item Chi phí vận hành hạ tầng Cache (Redis + Qdrant cluster): $\$1.500\text{/tháng}$
  \end{itemize}
  Hãy tính toán bằng tay:
  \begin{itemize}
    \item[a)] Tổng số tiền công ty phải trả cho việc gọi LLM mỗi tháng (30 ngày) khi \textbf{chưa có cache}.
    \item[b)] Tổng số tiền tiết kiệm ròng (\textit{Net Savings}) mỗi tháng sau khi trừ chi phí vận hành hạ tầng cache.
    \item[c)] Tỷ suất hoàn vốn đầu tư (\textit{Return on Investment - ROI}) của hạ tầng cache.
  \end{itemize}

  \item \textbf{Ước tính số lượng GPU vLLM cho hệ thống đọc tài liệu}:
  Hệ thống tiếp nhận tải đỉnh 500 QPS. Mỗi yêu cầu có độ dài prompt là $2.000$ tokens và độ dài output là $100$ tokens.
  Một máy chủ trang bị 8 GPU L40S có năng lực phục vụ tối đa $800$ decode tokens/giây và có bộ nhớ VRAM đủ duy trì tối đa $64$ requests đồng thời.
  Hãy tính toán số lượng máy chủ GPU tối thiểu cần chuẩn bị để không vượt quá $80\%$ công suất tải đỉnh.

  \item \textbf{Đánh đổi ngưỡng $\tau$ và tỷ lệ dương tính giả (False Positive Rate)}:
  Xét một hệ thống Semantic Cache có phân phối độ tương đồng giữa các câu hỏi cùng ý nghĩa $\mathcal{N}(\mu_1 = 0.93, \sigma_1 = 0.03)$ và giữa các câu hỏi khác ý nghĩa $\mathcal{N}(\mu_0 = 0.82, \sigma_0 = 0.04)$. Nếu đặt ngưỡng $\tau = 0.88$, hãy tính xác suất câu hỏi khác ý nghĩa bị nhận nhầm là trùng khớp (False Positive Hit Rate).

  \item \textbf{Định lượng tiết kiệm TFLOPS và TTFT nhờ Prefix-Aware Load Balancer}:
  Một cụm vLLM gồm 8 GPU H100 phục vụ ứng dụng lập trình viên Copilot.
  System prompt chứa tài liệu dự án dài $4.000$ tokens. Mỗi ngày có $100.000$ truy vấn chứa tiền tố này.
  Nếu dùng cân bằng tải Round-Robin ngây thơ: Tỷ lệ Cache Hit chỉ đạt $12.5\%$ (do phân bổ ngẫu nhiên đều cho 8 máy chủ).
  Nếu dùng Prefix-Aware Load Balancer: Tỷ lệ Cache Hit đạt $95\%$.
  Biết rằng pha Prefill cho $4.000$ tokens trên mô hình 70B tiêu tốn $560\text{ TFLOPs}$ và mất $400\text{ms}$.
  Hãy tính:
  \begin{itemize}
    \item[a)] Số lượng phép tính TFLOPs phần cứng GPU tiết kiệm được mỗi ngày.
    \item[b)] Tổng số giờ chờ đợi độ trễ TTFT được giải phóng cho người dùng mỗi ngày.
  \end{itemize}

  \item \textbf{Xác suất dịch chuyển khóa tiền tố trong Consistent Hashing}:
  Cụm vLLM đang hoạt động với $N = 10$ máy chủ GPU, quản lý $10.000$ mẫu tiền tố của các khách hàng doanh nghiệp.
  Do nhu cầu tăng trưởng, đội ngũ kỹ sư bổ sung thêm $2$ máy chủ GPU mới vào cụm ($N_{\text{new}} = 12$).
  Hãy tính:
  \begin{itemize}
    \item[a)] Số lượng mẫu tiền tố dự kiến bị dịch chuyển sang máy chủ mới theo nguyên lý Consistent Hashing.
    \item[b)] Số lượng mẫu tiền tố bị đảo lộn vị trí nếu sử dụng hàm băm chia dư truyền thống $\text{Hash}(K) \pmod N$.
  \end{itemize}

  \item \textbf{Tính Toán Ngân Sách Sai Số (Error Budget) \& Tốc Độ Đốt (Burn Rate)}:
  Một nền tảng AI doanh nghiệp cam kết chất lượng dịch vụ $\text{SLA} = 99.9\%$ cho khối lượng $N_{\text{total}} = 10.000.000$ truy vấn trong chu kỳ 30 ngày.
  \begin{itemize}
    \item[a)] Tính số lượng truy vấn lỗi tối đa được phép xảy ra trong cả tháng mà không bị vi phạm SLA ($N_{\text{allowed\_failures}}$).
    \item[b)] Nếu trong 1 giờ cao điểm đầu tiên, hệ thống tiếp nhận $50.000$ truy vấn và ghi nhận $720$ truy vấn bị lỗi (tương ứng tỷ lệ lỗi quan sát là $1.44\%$). Hãy tính tốc độ đốt ngân sách sai số ($\text{Burn Rate}$) trong giờ này.
    \item[c)] Với tốc độ lỗi trên, sau bao nhiêu giờ thì toàn bộ ngân sách sai số của cả tháng 30 ngày sẽ bị đốt sạch hoàn toàn?
  \end{itemize}

  \item \textbf{Tính Toán Tham Số Token Bucket \& Thời Gian Chờ Thử Lại (Retry-After)}:
  Một khách hàng doanh nghiệp thuộc gói Pro được cấp một xô thẻ Token Bucket với tốc độ nạp $R = 10\text{ tokens/giây}$ và dung lượng tối đa $B = 20\text{ tokens}$.
  Giả sử tại $t = 0.0\text{s}$, xô đang đầy nguyên vẹn $20$ tokens. Khách hàng gửi liên tiếp các đợt yêu cầu:
  - Đợt 1 (tại $t = 0.0\text{s}$): Yêu cầu $15$ tokens đồng thời.
  - Đợt 2 (tại $t = 0.0\text{s}$): Ngay lập tức yêu cầu thêm $10$ tokens nữa.
  - Đợt 3 (tại $t = 1.0\text{s}$): Sau 1 giây, yêu cầu $10$ tokens.
  Hãy tính số token còn lại trong xô, trạng thái chấp thuận/từ chối, và thời gian chờ thử lại $\Delta t_{\text{retry}}$ (nếu bị từ chối) cho từng đợt yêu cầu.
\end{enumerate}

\subsection{Nhóm C: Thực Hành Kỹ Nghệ \& Lập Trình (Implementation Labs)}

\begin{enumerate}[label=\textbf{C\arabic*.}]
  \item \textbf{Bộ cân bằng tải Least-Active-Connections có hỗ trợ Prefix Affinity}:
  Xây dựng lớp Python \code{ContextAwareLoadBalancer} quản lý danh sách $K$ máy chủ vLLM. Mỗi máy chủ ghi nhận số lượng kết nối đang mở (\code{active_conns}) và tập hợp các \code{system_prompt_hashes} đang được lưu trong cache VRAM của nó. Ưu tiên máy chủ có Prefix Cache Hit trước khi chọn máy chủ có tải thấp nhất.
  \item \textbf{Bộ điều phối Cascading Router có đo lường độ tự tin}:
  Triển khai lớp \code{CascadingModelRouter} gọi mô hình nhỏ trước; nếu điểm xác suất trung bình của các token sinh ra thấp hơn ngưỡng $\theta = 0.75$, tự động hủy kết quả và chuyển tiếp truy vấn lên mô hình lớn.
  \item \textbf{Token Bucket Rate Limiter phân tán đa người dùng}:
  Xây dựng bộ kiểm soát lưu lượng theo IP và API Key hỗ trợ burst capacity và nạp thẻ bất đồng bộ.
  \item \textbf{Bộ Định Tuyến GPU Đa Vùng Kèm Ngắt Mạch Tự Động (Multi-Region GPU Router)}:
  Cài đặt lớp Python \code{MultiRegionGPURouter} quản lý danh sách các vùng GPU phân tán (Singapore, Frankfurt, Virginia). Tự động tính toán độ trễ hiệu dụng $L_{\text{eff}} = L_{\text{base}} \times (1 + \frac{\text{Load}}{\text{Capacity}})$, tích hợp máy trạng thái ngắt mạch 3 pha (CLOSED, OPEN, HALF-OPEN), và tự động chuyển hướng sang nhà cung cấp đám mây thứ ba (\code{CLOUD_PROVIDER_FALLBACK}) khi toàn bộ cụm nội bộ gặp sự cố hoặc vượt ngưỡng SLA độ trễ.
  \item \textbf{Bộ Đo Ngân Sách Sai Số SLO \& Ước Tính Phân Vị Độ Trễ ($P50, P90, P99, P99.9$)}:
  Cài đặt lớp Python \code{SLOErrorBudgetCalculator} tính toán chính xác số lượng lỗi cho phép còn lại, tỷ lệ phần trăm ngân sách đã tiêu thụ, tốc độ đốt ngân sách Burn Rate, và thuật toán nội suy tuyến tính tính toán các phân vị độ trễ chuẩn mực $P50, P90, P99, P99.9$.
\end{enumerate}

\subsection{Nhóm D: Phân Tích Thất Bại \& Debugging Thực Tế (Post-Mortem Analysis)}

\begin{enumerate}[label=\textbf{D\arabic*.}]
  \item \textbf{Sự cố lộ dữ liệu số dư tài khoản ngân hàng qua Semantic Cache}:
  Một ngân hàng số triển khai Semantic Cache cho chatbot tư vấn tài chính.
  Một khách hàng VIP hỏi: ``Số dư tài khoản tiết kiệm của tôi là bao nhiêu?''. Chatbot gọi LLM và trả lời: ``Số dư của bạn là 500 triệu VNĐ''. Câu hỏi và câu trả lời này được lưu vào Semantic Cache toàn cục.
  Năm phút sau, một khách hàng phổ thông khác hỏi: ``Tôi muốn kiểm tra số dư tiết kiệm của mình''.
  Do độ tương đồng Cosine giữa hai câu hỏi đạt $0.94 > \tau = 0.88$, Semantic Cache lập tức trả về câu trả lời của khách VIP!
  Hãy viết báo cáo phân tích sự cố (Post-Mortem): Phân tích sai lầm cốt tử trong thiết kế Cache Key, và xây dựng lại kiến trúc \textbf{Tenant/User-Scoped Cache Key} an toàn tuyệt đối.
  \item \textbf{Sự cố sụp đổ dây chuyền do Cascading Storm}:
  Khi một mô hình SLM 8B bị lỗi định dạng phản hồi (trả về JSON sai cú pháp), toàn bộ $100\%$ lưu lượng bị dồn lên cụm LLM 70B phía sau. Cụm 70B bị quá tải khiến hàng đợi tràn, RAM cạn kiệt và sụp đổ hoàn toàn. Phân tích giải pháp Circuit Breaker và Shedding Load để cô lập sự cố.
\end{enumerate}

---

\section{Hướng Dẫn Lời Giải \& Bình Luận Kỹ Thuật Chi Tiết}
\label{sec:ch13-solutions}

\paragraph{Lời giải Câu A1 (Rò rỉ dữ liệu qua Semantic Cache)}:
Semantic Caching tìm kiếm sự tương đồng thuần túy dựa trên ngữ nghĩa của câu chữ mà mặc định không nhận biết được ngữ cảnh bảo mật của người hỏi.
Nếu hai người dùng thuộc hai công ty khác nhau (Tenant A và Tenant B) hoặc hai cá nhân khác nhau hỏi cùng một câu hỏi chung, việc dùng chung một kho vector cache toàn cục sẽ dẫn đến nguy cơ rò rỉ dữ liệu chéo nghiêm trọng.
\textbf{Giải pháp}: Bắt buộc phải áp dụng kiến trúc \textbf{Partitioned Cache Keys}:
Khóa tìm kiếm trong cache không chỉ là vector ngữ nghĩa mà phải đi kèm bộ lọc metadata bắt buộc:
\begin{equation*}
\text{Filter} = (\text{tenant\_id} == \text{current\_tenant}) \;\land\; \Big( \text{data\_classification} == \text{"PUBLIC"} \;\lor\; \text{user\_id} == \text{current\_user} \Big)
\end{equation*}
Chỉ những câu hỏi thuộc dữ liệu công khai dùng chung (như chính sách, tài liệu hướng dẫn) mới được chia sẻ giữa các người dùng.

\paragraph{Lời giải Câu A2 (Đánh đổi kích thước mô hình Embedding)}:
Trong trường hợp Cache Miss, người dùng phải chịu một hình phạt kép về độ trễ (\textit{Cache Miss Penalty}):
\begin{equation*}
\text{Total Latency}_{\text{miss}} = \text{Embedding Time} + \text{Vector Search Time} + \text{LLM Generation Time}
\end{equation*}
Nếu dùng mô hình embedding lớn mất $40\text{ms}$, mỗi lần trượt cache, hệ thống sẽ lãng phí thêm $40\text{ms}$ vô ích trước khi gọi được LLM.
Do đó, trong kiến trúc Semantic Cache, người ta luôn ưu tiên các mô hình embedding cực nhẹ (dưới 50MB, kích thước 384 chiều, chạy trên CPU/ONNX Runtime mất dưới $5\text{ms}$) để giảm thiểu tối đa hình phạt độ trễ khi Miss.

\paragraph{Lời giải Câu A3 (Chỉ số đánh giá TPOT)}:
Khi sử dụng giao thức Streaming (Server-Sent Events), người dùng nhìn thấy các từ xuất hiện liên tục ngay khi chúng được sinh ra.
Nếu TPOT đạt mức $40 - 50\text{ms/token}$ (tốc độ $20 - 25\text{ tokens/s}$), người dùng có cảm giác hệ thống phản hồi tức thì và đọc theo dòng văn bản một cách tự nhiên.
Dù toàn bộ câu trả lời 1000 tokens mất tổng cộng 40 giây (Total Latency lớn), người dùng không cảm thấy phải chờ đợi. Ngược lại, nếu không stream mà chờ đủ 40 giây mới hiển thị toàn bộ văn bản một lần, người dùng sẽ cho rằng hệ thống đã bị treo hoặc lỗi mạng!

\paragraph{Lời giải Câu A4 (Nguy cơ Cascading Storm)}:
Khi $90\%$ lưu lượng bị đẩy dồn sang mô hình 70B, cụm máy chủ lớn sẽ nhận tải gấp 3 lần năng lực thiết kế, dẫn đến hiện tượng hàng đợi phình to, VRAM cạn kiệt và sụp đổ dây chuyền (\textit{Cascading Failure}).
\textbf{Giải pháp}:
\begin{enumerate}
  \item Triển khai Circuit Breaker: Giới hạn lưu lượng fallback tối đa không vượt quá $30\%$ dung lượng cụm 70B.
  \item Shedding Load (Hạ tải có kiểm soát): Khi cụm 70B chạm ngưỡng $85\%$ VRAM hoặc hàng đợi $> 100$, hệ thống từ chối fallback và trả về câu trả lời của mô hình nhỏ kèm cảnh báo, hoặc gọi Cloud API dự phòng.
\end{enumerate}

\paragraph{Lời giải Câu A5 (Bảo toàn Cache của Consistent Hashing)}:
Với thuật toán băm chia dư ngây thơ $\text{Server} = \text{Hash}(K) \pmod N$: Khi một máy chủ GPU gặp sự cố sập ($N \to N-1$), mẫu số chia dư thay đổi khiến vị trí của hầu như toàn bộ $100\%$ các khóa tiền tố bị dịch chuyển sang máy chủ khác. Các GPU mới hoàn toàn chưa có KV Cache của các tiền tố này trong VRAM, dẫn đến $100\%$ Cache Miss và gây ra cơn bão Prefill Storm đánh sập toàn bộ các máy chủ còn lại trong cụm.
Ngược lại, với Consistent Hashing trên vòng $2^{32}$: Khi một máy chủ bị sập, chỉ có các dải khóa nằm trên cung tròn của máy chủ đó bị bàn giao cho máy chủ kế tiếp theo chiều kim đồng hồ (tỷ lệ $\frac{1}{N}$). Toàn bộ $\frac{N-1}{N}$ dải khóa còn lại trên vòng tròn vẫn giữ nguyên máy chủ phụ trách, bảo toàn nguyên vẹn bộ nhớ đệm KV Cache.

\paragraph{Lời giải Câu A6 (Ổn định kết nối Anycast BGP)}:
Anycast BGP hoạt động ở tầng Mạng (IP Layer). Do các biến động định tuyến trên Internet (Route Flapping), các gói tin tiếp theo của cùng một kết nối TCP có thể bị chuyển hướng sang một trung tâm dữ liệu khác. Vì máy chủ tại trung tâm dữ liệu mới không nắm giữ trạng thái kết nối TCP (TCP TCB), nó sẽ gửi cờ RST làm đứt gãy luồng Server-Sent Events (SSE) của người dùng.
\textbf{Giải pháp}:
\begin{enumerate}
  \item Chuyển sang sử dụng giao thức HTTP/3 trên nền QUIC: Giao thức QUIC sử dụng Connection ID độc lập với 4-tuple IP/Port, cho phép tiếp tục phiên truyền dữ liệu mà không bị ngắt quãng khi tuyến đường Anycast thay đổi.
  \item Triển khai tầng cân bằng tải nhất quán (Consistent Packet Forwarding như Maglev/Katran): Khi nhận gói tin của phiên đang mở tại PoP khác, Gateway tự động chuyển tiếp gói tin qua mạng cáp quang nội bộ về đúng PoP gốc.
\end{enumerate}

\paragraph{Lời giải Câu A7 (Cơ chế bùng nổ độ trễ đuôi $P99/P99.9$ dưới góc nhìn hàng đợi $M/M/c$)}:
Theo mô hình hàng đợi $M/M/c$ (hoặc xấp xỉ Kingman trong hàng đợi tổng quát $G/G/c$), thời gian chờ trung bình trong hàng đợi $W_q$ tỷ lệ thuận với hệ số sử dụng tài nguyên:
\begin{equation*}
W_q \approx \left(\frac{c_a^2 + c_s^2}{2}\right) \cdot \left(\frac{\rho}{1 - \rho}\right) \cdot \frac{1}{\mu}
\end{equation*}
Trong đó $\rho = \frac{\lambda}{c \mu}$ là hệ số tải GPU ($0 \le \rho < 1$).
Khi tải $\rho$ tăng từ $0.70$ lên $0.95$, hệ số tắc nghẽn $\frac{\rho}{1 - \rho}$ bùng nổ phi tuyến tính:
\begin{equation*}
\left. \frac{\rho}{1 - \rho} \right|_{\rho=0.70} = \frac{0.70}{0.30} \approx 2.33, \qquad \left. \frac{\rho}{1 - \rho} \right|_{\rho=0.95} = \frac{0.95}{0.05} = 19.0
\end{equation*}
Thời gian chờ trung bình tăng vọt $19.0 / 2.33 \approx 8.15$ lần!
Đặc biệt, phân phối thời gian chờ trong hàng đợi có đuôi hàm mũ:
\begin{equation*}
P(W_q > t) \approx C(c, \rho) e^{-c\mu(1-\rho)t} \implies t_p \propto \frac{1}{c\mu(1-\rho)} \ln \left( \frac{C(c, \rho)}{1 - p} \right)
\end{equation*}
Tại phân vị đuôi cực đoan $p = 0.99$ ($\ln \frac{1}{1-p} \approx 4.61$) và $p = 0.999$ ($\ln \frac{1}{1-p} \approx 6.91$), khi mẫu số $(1-\rho) \to 0$, độ trễ đuôi $P99$ và $P99.9$ phát nổ tiệm cận đứng. Mọi block VRAM đều bị chiếm giữ bởi các luồng decode dở dang, khiến các yêu cầu mới bị dồn ứ hàng chục giây tại hàng đợi Prefill.
\textbf{Lý do không đặt ngưỡng Autoscaling ở $90\%$}: Thời gian Cold-start của Pod GPU vLLM (kéo container, nạp weights mô hình 70B từ NVMe sang VRAM qua PCIe/NVLink) mất từ $45\text{s}$ đến $120\text{s}$. Nếu chờ tới $\rho = 0.90$ mới kích hoạt co giãn, bất kỳ đợt bùng nổ lưu lượng nào trong 60 giây chờ đợi cũng sẽ đẩy $\rho \to 1.0$, gây tràn hàng đợi bộ nhớ và sập dây chuyền toàn cụm. Ngưỡng an toàn chuẩn kỹ nghệ bắt buộc phải là $\rho = 70\% - 75\%$.

\paragraph{Lời giải Câu A8 (Chiến lược ngắt mạch Circuit Breaker và Fallback đa tầng)}:
Ba trạng thái hoạt động:
\begin{enumerate}
  \item \textbf{CLOSED (Đóng mạch)}: AI Gateway chuyển tiếp toàn bộ lưu lượng tới cụm GPU chính. Bộ đếm theo dõi số lỗi liên tiếp hoặc tỷ lệ lỗi trong cửa sổ trượt. Khi số lỗi liên tiếp chạm ngưỡng (ví dụ 5 lỗi) hoặc tỷ lệ lỗi vượt ngưỡng, mạch chuyển sang OPEN.
  \item \textbf{OPEN (Ngắt mạch)}: Cụm GPU được xác định là mất năng lực phục vụ. Gateway ngắn mạch lập tức (fail-fast), không gửi thêm request nào tới cụm này để tránh làm trầm trọng sự cố, đồng thời định tuyến tự động sang vùng thứ cấp hoặc nhà cung cấp đám mây dự phòng (\code{CLOUD_PROVIDER_FALLBACK}). Trạng thái duy trì trong thời gian hạ nhiệt $T_{\text{cooldown}}$ (30--60 giây).
  \item \textbf{HALF-OPEN (Nửa mở / Thăm dò)}: Sau thời gian $T_{\text{cooldown}}$, mạch cho phép một lượng nhỏ lưu lượng canary ($5\%$, hoặc 5--10 requests) đi qua thăm dò. Nếu toàn bộ lượt thăm dò thành công, mạch phục hồi về CLOSED. Nếu có bất kỳ lỗi nào xuất hiện, mạch bật lại OPEN và tăng gấp đôi $T_{\text{cooldown}}$ (exponential backoff).
\end{enumerate}
\textbf{Lý do bắt buộc thời gian chờ nguội và tỷ lệ $5\%$}: Khi cụm GPU vừa khởi động lại, bộ nhớ đệm KV Cache hoàn toàn nguội lạnh (Cold Cache). Nếu ngay lập tức mở $100\%$ tải (hàng ngàn QPS), hiện tượng bão yêu cầu Prefill (\textit{Thundering Herd}) sẽ lập tức ập tới, làm tràn VRAM và tái diễn sự cố sụp đổ ngay ở giây đầu tiên. Tỷ lệ canary $5\%$ giúp hệ thống khởi động mềm và làm ấm cache một cách an toàn.

\paragraph{Lời giải Câu A9 (Ý nghĩa toán học của Burn Rate trong quản trị SLO/SLA)}:
Tốc độ đốt ngân sách sai số được định nghĩa:
\begin{equation*}
\text{Burn Rate} = \frac{\text{Tỷ lệ lỗi quan sát được (Observed Error Rate)}}{1 - \text{SLA}}
\end{equation*}
Với $\text{SLA} = 99.9\%$, tỷ lệ lỗi tối đa được phép là $1 - 0.999 = 0.001$ ($0.1\%$).
Khi $\text{Burn Rate} = 1.0$, hệ thống tiêu thụ vừa vặn $100\%$ ngân sách trong toàn bộ chu kỳ 30 ngày (tốc độ bền vững).
\textbf{Lý do giám sát đa cửa sổ (1h và 6h)}: Nếu chỉ đánh giá sau 30 ngày thì mọi khắc phục đều quá trễ. Cửa sổ 1 giờ giúp bắt các thảm họa bùng phát đột ngột, trong khi cửa sổ 6 giờ phát hiện các lỗi rò rỉ âm ỉ kéo dài mà cửa sổ ngắn dễ bỏ sót.
\textbf{Lý do chọn $\text{Burn Rate} = 14.4\times$ làm ngưỡng Critical Pager}:
Trong chu kỳ 30 ngày có $30 \times 24 = 720$ giờ. Tỷ lệ phần trăm ngân sách bị đốt trong 1 giờ là:
\begin{equation*}
\% \text{Budget Consumed} = \text{Burn Rate} \times \frac{1\text{ giờ}}{720\text{ giờ}} = \frac{14.4}{720} = 0.02 = \mathbf{2\%}
\end{equation*}
Theo quy chuẩn Google SRE, sự cố nào thiêu rụi từ $2\%$ toàn bộ ngân sách sai số của cả tháng chỉ trong vòng 1 giờ đồng hồ là thảm họa nghiêm trọng, vì toàn bộ ngân sách tháng sẽ bị xóa sổ hoàn toàn sau $720 / 14.4 = 50$ giờ (hơn 2 ngày). Do đó, hệ thống phải kích hoạt Pager gọi kỹ sư trực ca can thiệp khẩn cấp.

\paragraph{Lời giải Bài B1 (Tính toán ROI Semantic Cache)}:
\begin{itemize}
  \item[a)] Tổng số truy vấn mỗi tháng:
  \begin{equation*}
  N_{\text{month}} = 5.000.000 \text{ req/ngày} \times 30 \text{ ngày} = 150.000.000 \text{ requests}
  \end{equation*}
  Chi phí khi chưa có cache:
  \begin{equation*}
  \text{Cost}_{\text{no\_cache}} = 150.000.000 \times \$0.01 = \mathbf{\$1.500.000\text{/tháng}}
  \end{equation*}

  \item[b)] Hiệu quả giảm tải của bộ nhớ đệm hai tầng:
  Tổng tỷ lệ Hit Rate kết hợp:
  \begin{equation*}
  \text{HitRate}_{\text{total}} = 15\% (\text{Exact}) + 25\% (\text{Semantic}) = 40\%
  \end{equation*}
  Số lượng truy vấn được giải đáp trực tiếp từ cache mà không cần gọi LLM:
  \begin{equation*}
  N_{\text{saved}} = 150.000.000 \times 40\% = 60.000.000 \text{ requests}
  \end{equation*}
  Số tiền tiết kiệm thô từ việc không gọi LLM:
  \begin{equation*}
  \text{Gross Savings} = 60.000.000 \times \$0.01 = \$600.000\text{/tháng}
  \end{equation*}
  Số tiền tiết kiệm ròng sau khi trừ chi phí hạ tầng cache:
  \begin{equation*}
  \text{Net Savings} = \$600.000 - \$1.500 = \mathbf{\$598.500\text{/tháng}}
  \end{equation*}

  \item[c)] Tỷ suất hoàn vốn đầu tư (ROI):
  \begin{equation*}
  \text{ROI} = \frac{\$598.500}{\$1.500} \times 100\% = \mathbf{39.900\%}!
  \end{equation*}
  Khoản đầu tư vào Semantic Cache mang lại lợi nhuận gấp gần 400 lần chi phí bỏ ra!
\end{itemize}

\paragraph{Lời giải Bài B2 (Ước tính số lượng GPU L40S)}:
Tải đỉnh yêu cầu sinh token:
\begin{equation*}
\text{Decode Rate} = 500 \text{ QPS} \times 100 \text{ tokens} = 50.000 \text{ tokens/giây}
\end{equation*}
Mỗi máy chủ xử lý tối đa 800 tokens/giây. Để giữ tải ở mức an toàn $\le 80\%$ công suất:
\begin{equation*}
\text{Capacity per server (80\%)} = 800 \times 0.8 = 640 \text{ tokens/giây}
\end{equation*}
Số lượng máy chủ cần thiết theo thông lượng sinh token:
\begin{equation*}
N_{\text{servers, throughput}} = \frac{50.000}{640} = 78.125 \implies 79 \text{ máy chủ}
\end{equation*}
Kiểm tra ràng buộc về số lượng kết nối đồng thời:
Giả sử thời gian trung bình của mỗi request là $T = \text{TTFT} + \text{DecodeTime} = 0.5\text{s} + \frac{100}{30}\text{s} \approx 3.83\text{s}$.
Theo định luật Little ($L = \lambda W$), số request đồng thời vào giờ cao điểm:
\begin{equation*}
L = 500 \times 3.83 \approx 1.915 \text{ requests}
\end{equation*}
Số máy chủ cần thiết theo dung lượng kết nối (mỗi máy tối đa $64 \times 0.8 = 51.2$ requests):
\begin{equation*}
N_{\text{servers, concurrency}} = \frac{1.915}{51.2} \approx 37.4 \implies 38 \text{ máy chủ}
\end{equation*}
Lấy giá trị lớn nhất giữa hai ràng buộc: Hệ thống cần tối thiểu \textbf{79 máy chủ GPU L40S} (tương đương $79 \times 8 = 632$ GPU) để đảm bảo an toàn tuyệt đối ở tải đỉnh!

\paragraph{Lời giải Bài B3 (Đánh đổi ngưỡng $\tau$ và False Hit Rate)}:
Độ tương đồng giữa hai câu hỏi khác ý nghĩa tuân theo $\mathcal{N}(\mu_0 = 0.82, \sigma_0 = 0.04)$.
Xác suất để một câu hỏi khác ý nghĩa có độ tương đồng $\ge \tau = 0.88$ (False Positive):
\begin{equation*}
Z = \frac{\tau - \mu_0}{\sigma_0} = \frac{0.88 - 0.82}{0.04} = 1.5
\end{equation*}
Tra bảng phân phối chuẩn tắc $P(Z \ge 1.5) \approx 0.0668 = 6.68\%$.
Con số này cho thấy nếu chỉ dựa vào Cosine similarity với ngưỡng $0.88$, cứ 100 câu hỏi khác ý nghĩa sẽ có gần 7 câu bị trả lời sai! Do đó, việc kết hợp kiểm định thực thể (Entity Extractor) là bắt buộc.

\paragraph{Lời giải Bài B4 (Định lượng tiết kiệm TFLOPS và TTFT)}:
\begin{itemize}
  \item[a)] Số lượng truy vấn được hưởng Prefix Cache Hit mỗi ngày:
  \begin{itemize}
    \item Khi dùng Round-Robin: $100.000 \times 12.5\% = 12.500$ lượt $\implies 87.500$ lượt Cache Miss.
    \item Khi dùng Prefix-Aware: $100.000 \times 95\% = 95.000$ lượt $\implies 5.000$ lượt Cache Miss.
  \end{itemize}
  Số lượt tính toán Prefill được cắt giảm mỗi ngày:
  \begin{equation*}
  \Delta N = 87.500 - 5.000 = 82.500 \text{ requests}
  \end{equation*}
  Tổng số TFLOPs tiết kiệm được mỗi ngày:
  \begin{equation*}
  \text{TFLOPs Saved} = 82.500 \times 560\text{ TFLOPs} = \mathbf{46.200.000\text{ TFLOPs}} = \mathbf{46.2\text{ ExaFLOPs/ngày}}
  \end{equation*}

  \item[b)] Tổng số giờ chờ đợi TTFT được giải phóng cho người dùng:
  \begin{equation*}
  \text{Latency Hours Saved} = \frac{82.500 \times 0.400\text{ giây}}{3.600\text{ giây/giờ}} = \frac{33.000}{3.600} \approx \mathbf{9.17\text{ giờ/ngày}}
  \end{equation*}
  Người dùng tiết kiệm được hơn 9 giờ đồng hồ chờ đợi phản hồi mỗi ngày!
\end{itemize}

\paragraph{Lời giải Bài B5 (Xác suất dịch chuyển khóa trong Consistent Hashing)}:
\begin{itemize}
  \item[a)] Với thuật toán Consistent Hashing:
  Khi tăng số máy chủ từ $N_{\text{old}} = 10$ lên $N_{\text{new}} = 12$, tỷ lệ khóa tiền tố bị dịch chuyển sang hai máy chủ mới là:
  \begin{equation*}
  \text{Tỷ lệ dịch chuyển} = \frac{N_{\text{new}} - N_{\text{old}}}{N_{\text{new}}} = \frac{12 - 10}{12} = \frac{2}{12} \approx 16.67\%
  \end{equation*}
  Số lượng mẫu tiền tố bị di dời:
  \begin{equation*}
  N_{\text{moved, consistent}} = 10.000 \times 16.67\% \approx \mathbf{1.667\text{ mẫu}}
  \end{equation*}
  Hơn $83.33\%$ ($8.333$ mẫu) vẫn giữ nguyên vị trí trên các GPU cũ, duy trì tỷ lệ Cache Hit ổn định.

  \item[b)] Với thuật toán băm chia dư $\text{Hash}(K) \pmod N$:
  Tỷ lệ các khóa giữ nguyên vị trí là tỷ lệ thỏa mãn $\text{Hash}(K) \pmod{10} = \text{Hash}(K) \pmod{12}$, tức $\text{Hash}(K) \pmod{\text{gcd}(10, 12)} = \text{Hash}(K) \pmod 2$.
  Tỷ lệ giữ nguyên vị trí chỉ là:
  \begin{equation*}
  \frac{\text{gcd}(10, 12)}{\text{lcm}(10, 12)} = \frac{2}{60} \approx 3.33\%
  \end{equation*}
  Số lượng mẫu tiền tố bị đảo lộn vị trí:
  \begin{equation*}
  N_{\text{moved, modulo}} = 10.000 \times (1 - 0.0333) = \mathbf{9.667\text{ mẫu}} \quad (\text{Mất trắng hơn } 96.67\% \text{ cache!})
  \end{equation*}
\end{itemize}

\paragraph{Lời giải Bài B6 (Tính toán ngân sách sai số Error Budget và Burn Rate)}:
\begin{itemize}
  \item[a)] Số lượng truy vấn lỗi tối đa được phép xảy ra trong 30 ngày:
  \begin{equation*}
  N_{\text{allowed\_failures}} = N_{\text{total}} \times (1 - \text{SLA}) = 10.000.000 \times (1 - 0.999) = 10.000.000 \times 0.001 = \mathbf{10.000\text{ truy vấn}}
  \end{equation*}

  \item[b)] Tỷ lệ lỗi quan sát được trong 1 giờ cao điểm:
  \begin{equation*}
  \text{Error Rate}_{\text{1h}} = \frac{720}{50.000} = 0.0144 = 1.44\%
  \end{equation*}
  Tốc độ đốt ngân sách sai số (Burn Rate):
  \begin{equation*}
  \text{Burn Rate} = \frac{\text{Error Rate}_{\text{1h}}}{1 - \text{SLA}} = \frac{0.0144}{0.001} = \mathbf{14.4\times}
  \end{equation*}

  \item[c)] Thời gian toàn bộ ngân sách sai số của cả tháng bị tiêu sạch:
  \begin{equation*}
  \text{Time to Exhaust} = \frac{30 \times 24\text{ giờ}}{\text{Burn Rate}} = \frac{720}{14.4} = \mathbf{50.0\text{ giờ}}
  \end{equation*}
  Chỉ sau hơn 2 ngày với tốc độ lỗi này, công ty sẽ tiêu sạch toàn bộ hạn mức sai số cam kết cho cả tháng!
\end{itemize}

\paragraph{Lời giải Bài B7 (Tính toán tham số Token Bucket và Retry-After)}:
Các tham số: Dung lượng tối đa $B = 20\text{ tokens}$, tốc độ nạp $R = 10\text{ tokens/giây}$. Tại $t = 0.0\text{s}$, xô đầy: $\text{tokens}(0.0) = 20.0$.
\begin{itemize}
  \item \textbf{Đợt 1 (tại $t = 0.0\text{s}$, yêu cầu $15$ tokens)}:
  Do $20.0 \ge 15$, yêu cầu được \textbf{CHẤP THUẬN}.
  Số token còn lại trong xô: $20.0 - 15 = \mathbf{5.0\text{ tokens}}$.
  Thời gian chờ thử lại: $\Delta t_{\text{retry}} = \mathbf{0.0\text{s}}$.

  \item \textbf{Đợt 2 (ngay tại $t = 0.0\text{s}$, yêu cầu thêm $10$ tokens)}:
  Thời gian trôi qua $\Delta t = 0.0\text{s} \implies$ chưa nạp thêm token nào.
  Số token hiện có là $5.0 < 10$. Yêu cầu bị \textbf{TỪ CHỐI} (mã phản hồi HTTP 429 Too Many Requests).
  Số token thiếu hụt: $\Delta_{\text{deficit}} = 10 - 5.0 = 5.0\text{ tokens}$.
  Thời gian chờ thử lại khuyến nghị trong tiêu đề \code{Retry-After}:
  \begin{equation*}
  \Delta t_{\text{retry}} = \frac{\Delta_{\text{deficit}}}{R} = \frac{5.0}{10} = \mathbf{0.5\text{ giây}}
  \end{equation*}
  Số token trong xô được bảo toàn nguyên vẹn: $\mathbf{5.0\text{ tokens}}$.

  \item \textbf{Đợt 3 (tại $t = 1.0\text{s}$, yêu cầu $10$ tokens)}:
  Thời gian trôi qua từ lần cập nhật trước: $\Delta t = 1.0 - 0.0 = 1.0\text{s}$.
  Số token nạp mới: $\Delta t \times R = 1.0 \times 10 = 10.0\text{ tokens}$.
  Số token trước khi tiêu thụ: $\min(B, 5.0 + 10.0) = \min(20.0, 15.0) = 15.0\text{ tokens}$.
  Do $15.0 \ge 10$, yêu cầu được \textbf{CHẤP THUẬN}.
  Số token còn lại trong xô: $15.0 - 10 = \mathbf{5.0\text{ tokens}}$.
  Thời gian chờ thử lại: $\Delta t_{\text{retry}} = \mathbf{0.0\text{s}}$.
\end{itemize}

\paragraph{Lời giải Bài C1, C2, C3, C4, C5 (Hiện thực hóa hệ thống)}:
Toàn bộ mã nguồn sản xuất hoàn chỉnh, kiểm thử tự động với $100\%$ tỷ lệ bao phủ, và tài liệu kỹ thuật chi tiết cho các bài tập thực hành được cung cấp trực tiếp trong tệp kiểm thử chuẩn của dự án:
\begin{itemize}
  \item \textbf{Bài C1} (\code{ContextAwareLoadBalancer}): Cân bằng tải có tính toán Prefix Cache Affinity.
  \item \textbf{Bài C2} (\code{CascadingModelRouter}): Điều phối mô hình phân cấp dựa trên ngưỡng tin cậy token.
  \item \textbf{Bài C3} (\code{DistributedTokenBucketRateLimiter}): Kiểm soát lưu lượng theo tầng gói cước (Free, Pro, Enterprise) với nạp thẻ vi sai thời gian.
  \item \textbf{Bài C4} (\code{MultiRegionGPURouter}): Định tuyến GPU đa vùng địa lý, bù trễ tải động và máy trạng thái ngắt mạch 3 pha (CLOSED, OPEN, HALF-OPEN).
  \item \textbf{Bài C5} (\code{SLOErrorBudgetCalculator}): Tính toán ngân sách sai số, tốc độ đốt Burn Rate và phân vị độ trễ ($P50, P90, P99, P99.9$) theo phương pháp nội suy tuyến tính.
\end{itemize}
Chi tiết xem tại: \code{tests/test_ch13_system_design.py}.

\paragraph{Lời giải Bài D1 (Báo cáo sự cố lộ số dư và kiến trúc khắc phục)}:
\begin{enumerate}
  \item \textbf{Nguyên nhân cốt tử}: Toàn bộ các câu hỏi đều được đưa vào cùng một không gian vector chung mà không phân loại xem truy vấn đó có chứa dữ liệu cá nhân hay không.
  \item \textbf{Kiến trúc khắc phục}:
    \begin{itemize}
      \item Gắn nhãn Cacheability: Chỉ cho phép ghi vào Cache những truy vấn được gắn nhãn \code{CACHEABLE=TRUE} (chính sách, lãi suất chung). Các câu hỏi số dư bắt buộc gắn cờ \code{CACHE-CONTROL: NO-STORE}.
      \item Cấu trúc Compound Cache Key: Khóa cache phải băm kết hợp:
      \begin{equation*}
      \text{Key} = \text{Hash}(\text{tenant\_id} + \text{user\_id} + \text{prompt})
      \end{equation*}
    \end{itemize}
\end{enumerate}

\paragraph{Lời giải Bài D2 (Khắc phục Cascading Storm)}:
Triển khai mẫu Circuit Breaker: Khi tỷ lệ lỗi hoặc thời gian đáp ứng của cụm 70B vượt quá ngưỡng P99 ($> 3\text{s}$), Gateway tạm thời ngắt mạch fallback trong 10 giây và trả về câu trả lời tốt nhất có thể của mô hình 8B hoặc thông báo bảo trì tạm thời.
